\documentclass[11pt]{article}

\usepackage[utf8]{inputenc}
\usepackage[T1]{fontenc}
\usepackage{amsmath, amssymb, amsthm}
\usepackage{mathtools}
\usepackage[margin=1in]{geometry}
\usepackage{hyperref}
\usepackage{graphicx}



\theoremstyle{plain}
\newtheorem*{theorem*}{Theorem}
\newtheorem*{lemma*}{Lemma}
\newtheorem*{corollary*}{Corollary}

\theoremstyle{definition}
\newtheorem*{definition}{Definition}
\newtheorem*{example}{Example}
\newtheorem*{remark}{Remark}
\newtheorem*{remarks}{Remarks}
\newtheorem*{notation}{Notation}
\newtheorem*{assertion}{Assertion}
\newtheorem*{claim}{Claim}


\newcommand{\R}{\mathbb{R}}
\newcommand{\C}{\mathbb{C}}
\newcommand{\Z}{\mathbb{Z}}
\newcommand{\Q}{\mathbb{Q}}
\newcommand{\nbd}{neighborhood}
\DeclareMathOperator{\ord}{ord}
\DeclareMathOperator{\res}{res}
\DeclareMathOperator{\discr}{discr}
\DeclareMathOperator{\subres}{subres}
\DeclareMathOperator{\subdiscr}{subdiscr}
\DeclareMathOperator{\coeff}{coeff}
\DeclareMathOperator{\cont}{cont}
\DeclareMathOperator{\prm}{prim}
\DeclareMathOperator{\red}{red}
\DeclareMathOperator{\der}{der}
\DeclareMathOperator{\gsfd}{gsfd}
\DeclareMathOperator{\ldcf}{ldcf}
\DeclareMathOperator{\ldt}{ldt}
\DeclareMathOperator{\gcdop}{gcd}
\DeclareMathOperator{\PROJ}{PROJ}
\DeclareMathOperator{\APROJ}{APROJ}



\title{An Improved Projection Operation\\for Cylindrical Algebraic Decomposition\\[1em]\large Chapters 2, 3 and 4}
\author{Scott McCallum}
\date{February 1985}

\begin{document}

\maketitle

\section*{Chapter Two: Mathematical Preliminaries}

This chapter presents mathematical background material which the reader
should find helpful in reading Chapters 3 and 4.

Our discussion of projection of algebraic varieties presented in Chapter 3
involves a rather careful study of the properties of algebraic sets in the
neighborhood of a particular point (such properties are called
\emph{local properties}). A basic tool in local analysis of algebraic sets
is the (multiple) power series. Functions having power series expansions,
or analytic functions, thus enter the discussion in a natural way.

There are a number of texts which include a discussion of the elementary
properties of analytic functions of several complex variables (for
example, [GRO65], [BMA48], and [KAP66]). Analytic functions of real
variables are best studied with the aid of complex variables: Chapter 2 of
[BMA48] contains a comparison between the real and complex cases.
Section 1 of this chapter is a collection of many basic results about
analytic functions (of both real and complex variables). For proofs one is
generally referred to texts.

The original cylindrical algebraic decomposition (cad) algorithm [COL75]
decomposes $\R^n$ into semi-algebraic subsets which Collins called
\emph{cells}. It has been subsequently observed [KAH78] that the cells
produced by this decomposition of $n$-space are actually bona fide cells
in the sense of topology: that is, each cell is homeomorphic to the open
unit ball in $\R^i$, for some $i$, $0 \le i \le n$. What is further true
is that each cell is homeomorphic to an open unit ball via a mapping which
is \emph{analytic}: thus each cell is an analytic $i$-dimensional
submanifold of $\R^n$, for some $i$. This smoothness property of the cells
turns out to be quite important in developing an improved projection
operation for the cad algorithm.

Section 2 develops the concept of submanifold as far as needed for our
purposes. Although the material is standard, it does not appear in quite
the form we require in any of the texts. Our presentation is tailored to
the needs of our later chapters.

Section 3 is a collection of miscellaneous results which will prove useful
in subsequent chapters.

\section{Analytic Functions of Several Variables}

We assume that the reader is familiar with the elementary theory of
analytic functions of a single complex variable. We discuss the notion of
a multiple power series and that of an analytic function of several (real
or complex) variables. We present statements of many basic theorems about
such functions, often referring the reader to standard texts for the
proofs.

We present a short summary of the material presented in this section. A
function $f(x_1,\ldots,x_n)$ (the $x_i$ real or complex variables) is said
to be analytic if it has a (multiple) power series representation about
each point of its domain. An analytic function of complex variables is
also termed \emph{holomorphic}. An analytic function is continuous and
has continuous partial derivatives of all orders. A function defined as
the sum of a convergent power series is analytic, and its partial
derivatives can be obtained by differentiating the defining series
term-by-term. Sums, products and quotients (assuming nonzero denominator)
of analytic functions are analytic.

Let $\R$ denote the field of real numbers, and let $\C$ denote the field
of complex numbers, i.e.\ numbers of the form $z = x + iy$, where $x$ and
$y$ are elements of $\R$ and $i$ is a square root of $-1$. \emph{Throughout
this section $K$ will denote either $\R$ or $\C$}. $K^n$ will denote the
Cartesian product $K \times \cdots \times K$ of $n$ copies of $K$. In this
section, unless otherwise specified, we will use the notation
$x = (x_1, \ldots, x_n)$ for points of $K^n$.

\begin{definition}
Let $c = (c_1, \ldots, c_n)$ be a point of $K^n$. A \emph{power series
about $c$ over $K$} is an expression of the form
\begin{equation}
\sum_{i_1, \ldots, i_n = 0}^{\infty}
a_{i_1, \ldots, i_n} (x_1 - c_1)^{i_1} \cdots (x_n - c_n)^{i_n}
\tag{2.1.1}
\end{equation}
where the coefficients $a_{i_1, \ldots, i_n}$ are elements of $K$.
\end{definition}

\begin{remark}
If $n > 1$ then the power series (2.1.1) is a multiple power series
because of the multiple index $i_1, \ldots, i_n$. We can, however,
arrange the terms to form a simple series, for example, with $n = 2$:
\[
\begin{aligned}
a_{00} &+ a_{10}(x_1 - c_1) + a_{01}(x_2 - c_2) + a_{20}(x_1 - c_1)^2
+ a_{11}(x_1 - c_1)(x_2 - c_2) \\
&+ a_{02}(x_2 - c_2)^2 + a_{30}(x_1 - c_1)^3 + \cdots .
\end{aligned}
\]
We are thus, in this example, sweeping out all the combinations
$(i_1, i_2)$ by following diagonals of the corresponding array (see
fig.~2.1.1).
\end{remark}

\begin{center}
\textit{[Figure 2.1.1: diagram showing diagonal enumeration of pairs
$(i_1, i_2)$ in a two-dimensional grid with $i_1$ labelling columns
$0,1,2,3,4$ and $i_2$ labelling rows $0,1,2,3,4$, with arrows indicating
the diagonal sweep.]}
\end{center}

Another arrangement of the terms is obtained by going around the sides of
squares of increasing size, that is, choosing the successive pairs
$(i_1, i_2)$ as follows:
\[
(0,0),\ (1,0),\ (1,1),\ (0,1),\ (2,0),\ (2,1),\ (2,2),\ (1,2),\ (0,2),\
(3,0),\ \ldots .
\]

If $x = (x_1, \ldots, x_n)$ is a point of $K^n$ then one can ask whether
the series (2.1.1) has some arrangement as a convergent simple series. If
there exists an arrangement of the series as a simple series for which
one has absolute convergence, then by theorem 28 on p.~333 of [KAP52],
the series is absolutely convergent for every arrangement as a simple
series, and the sum is the same for all arrangements. We shall say,
simply, that the series (2.1.1) is absolutely convergent if it is
absolutely convergent for some arrangement as a simple series.

If $c \in K^n$ then a \emph{neighborhood} of $c$ is an open subset $W$ of
$K^n$ containing $c$. Let $c \in K^n$ and let
$r = (r_1, \ldots, r_n) \in \R^n$, with $r_i > 0$ for $1 \le i \le n$.
Consider the neighborhood
\[
\{ x \in K^n : |x_i - c_i| < r_i,\ 1 \le i \le n \}
\]
of $c$. In the case $K = \R$, we call this neighborhood a \emph{box} about
$c$ and denote it by $B(c;r)$. In case $K = \C$, we call this neighborhood
the \emph{polydisc} about $c$ of polyradius $r$ and denote it by
$\Delta(c;r)$. Note that if $c \in \R^n$, then
$B(c;r) = \Delta(c;r) \cap \R^n$. Let $1 \le s \le n$ and let
$\pi: K^n \to K^s$ be the projection
$\pi(x_1, \ldots, x_n) = (x_1, \ldots, x_s)$. Let $K = \R$ and let $B$ be
the box $B(c;r)$ in $\R^n$. Then we shall denote by $B^{(s)}$ the image
of $B$ under $\pi$, i.e.
\[
B^{(s)} = \{ (x_1, \ldots, x_s) \in \R^s : |x_i - c_i| < r_i,\
1 \le i \le s \}.
\]
Let $K = \C$ and let $\Delta$ be the polydisc $\Delta(c;r)$ in $\C^n$. We
similarly denote by $\Delta^{(s)}$ the image of $\Delta$ under $\pi$.

\begin{definition}
Let $U$ be an open subset of $K^n$ and let $f: U \to K$ be a function.
Then $f$ is said to be \emph{analytic in $U$} if each point $c$ of $U$
has a neighborhood $W \subseteq U$ such that $f$ has a power series
about $c$ over $K$
\begin{equation}
f(x_1, \ldots, x_n) = \sum_{i_1, \ldots, i_n = 0}^{\infty}
a_{i_1, \ldots, i_n} (x_1 - c_1)^{i_1} \cdots (x_n - c_n)^{i_n}
\tag{2.1.3}
\end{equation}
which is absolutely convergent for every point $x$ of $W$.
\end{definition}

\begin{example}
Every polynomial in $x_1, \ldots, x_n$ over $K$ is analytic in the whole
of $K^n$.
\end{example}

\begin{remark}
If $K = \C$ and $U \subseteq K^n$ is open, then a function $f: U \to K$
analytic in $U$ is also termed \emph{holomorphic in $U$}.
\end{remark}

\begin{remark}
In the case $K = \R$, the above is the definition of analytic function
that appears in most texts. In the case $K = \C$, one sometimes finds an
alternative, equivalent definition of holomorphic function: a
complex-valued function $f(z_1, \ldots, z_n)$ defined in the open subset
$U$ of $\C^n$ is said to be holomorphic in $U$ if it is continuous in $U$
and has continuous partial derivatives $\dfrac{\partial f}{\partial z_i}$
in $U$. That this definition is equivalent to our definition for the
case $K = \C$ follows from Theorems 2.1.1 and 2.1.8 below.
\end{remark}

\begin{theorem*}[2.1.1]
Let $U \subseteq K^n$ be open and let $f: U \to K$ be an analytic
function. Then $f$ is continuous and has continuous partial derivatives
of all orders, given by:
\[
\frac{\partial^{i_1 + \cdots + i_n} f}
{\partial x_1^{i_1} \cdots \partial x_n^{i_n}}(c)
= (i_1!) \cdots (i_n!)\, a_{i_1, \ldots, i_n}
\]
where the $a_{i_1, \ldots, i_n}$ are the coefficients of the power series
expansion~(2.1.3) of $f$ about the point $c$ of $U$.
\end{theorem*}

\textit{References for proof.} For the real case, (i.e.\ $K = \R$), we
refer the reader to Sec.~6-20 of [KAP52]; and for the complex case,
($K = \C$), to Theorem 1 of Chapter II of [BMA48]. $\square$

\textit{Remark 1:} It follows from Theorem 2.1.1 that the power series
expansion of an analytic function about a point is unique.

\textit{Remark 2:} The converse of Theorem 2.1.1 holds in the case
$K = \C$ (this is Theorem 2.1.8) but not in the case $K = \R$, as is
shown by the example $f(x) = e^{-1/x^2}$ (discussed in Sec.~6-17 of
[KAP52]).

A function defined as the sum of a convergent power series is, as might
be expected, analytic:

\begin{theorem*}[2.1.2]
Let $c = (c_1, \ldots, c_n)$ be a point of $K^n$ and let the power
series over $K$
\begin{equation}
\sum_{i_1, \ldots, i_n = 0}^{\infty}
a_{i_1, \ldots, i_n} (x_1 - c_1)^{i_1} \cdots (x_n - c_n)^{i_n}
\tag{2.1.4}
\end{equation}
converge (for some arrangement as a simple series) for
$x_1 = c_1 + r_1,\ \ldots,\ x_n = c_n + r_n$, where each $r_i > 0$. Let
\[
D = \{ x = (x_1, \ldots, x_n) \in K^n : |x_i - c_i| < r_i,\
1 \le i \le n \}.
\]
Then the series~(2.1.4) is absolutely convergent for every $x \in D$, and
the sum of the series, say $f(x)$, is an analytic function in $D$.
Moreover, every partial derivative (of any order) of $f$ is analytic in
$D$, and its power series expansion about $c$, absolutely convergent for
every $x \in D$, is obtained by differentiating~(2.1.4) term-by-term.
\end{theorem*}

\begin{proof}
We first present a proof for the case $K = \C$. For the complex case let
us imagine that $z_i$ replaces each occurrence of $x_i$ in the statement
of the theorem. Note that in this case, the neighborhood $D$ of $c$ is
the polydisc $\Delta(c;r)$ about $c$. The first part of the conclusions,
as to the absolute convergence of the series~(2.1.4), follows from the
$n$-variable analogue of Theorem 54 in [KAP66]. By analogy with
Theorem 55 of [KAP66], let $E$ be the set of points $z \in \C^n$ for
which the series~(2.1.4) converges and let $E^i$ be the interior of $E$.
As the series~(2.1.4) converges for every $z \in D$, we have $D \subseteq E$.
Consequently, $D \subseteq E^i$. Hence the hypothesis of the $n$-variable
analogue of Theorem 55 of [KAP66] is satisfied. The holomorphicity of
the function $f$ now follows by the $n$-variable analogue of Theorem 56
of [KAP66]. The proof of this $n$-variable analogue of Theorem 56
yields the required results about the partial derivatives of $f$. The
theorem is now proved for the case $K = \C$.

Let us now deal with the real case. In this case the neighborhood $D$ of
$c$ is the box $B(c;r)$ about $c$. The first part of the conclusions as
to the convergence of~(2.1.4) follows from the $n$-variable analogue of
Theorem 54 in [KAP66]. By the complex case, the series~(2.1.4) (with the
$x_i$'s replaced by $z_i$'s) is absolutely convergent in
$\Delta_c := \Delta(c;r)$ and its sum, say $F(z)$, is holomorphic in
$\Delta_c$. Moreover, every partial derivative of $F$ is holomorphic in
$\Delta_c$, and its power series expansion about $c$ is obtained by
differentiating~(2.1.4) term-by-term. It is not hard to see that, for
every point $d$ of $D$, the power series expansion of $F(z)$ about $d$
has real coefficients. Hence, when $x_i$'s are substituted for $z_i$'s in
this power series about $d$, a power series representation for $f(x)$ is
obtained. It follows that $f$ is analytic in $D$.

The proof of the assertion concerning the partial derivatives of $f$ is
straightforward.
\end{proof}

It is well-known that sums, products and quotients of analytic functions
are analytic. We state this as

\begin{theorem*}[2.1.3]
Let $U$ be an open subset of $K^n$ and let $f$ and $g$ be analytic in
$U$. Then the functions $f + g$ and $fg$ (defined pointwise) are analytic
in $U$. If $f \ne 0$ in $U$ then the function $1/f$ is analytic in $U$.
\end{theorem*}

\begin{example}
Any rational function $f(x_1, \ldots, x_n)/g(x_1, \ldots, x_n)$ (where
$f$ and $g$ are polynomials over $K$) is analytic in the region of $K^n$
where $g \ne 0$.
\end{example}

\begin{theorem*}[2.1.4 (Identity Theorem)]
Let $D \subseteq K^n$ be a connected open set and let $f$ be analytic in
$D$. If $f(x) = 0$ for all points $x$ in a nonempty open subset $U$ of
$D$, then $f \equiv 0$ everywhere in $D$.
\end{theorem*}

\textit{References for proof.} The complex identity theorem is Theorem 6
of Ch.~I, [GRO65]. The real identity theorem is the theorem on page 53
of [JOH75]. $\square$

Let $U$ be an open subset of $K^n$ and let $G: U \to K^m$ be a mapping.
Then there are $m$ functions $g_1, \ldots, g_m$ from $U$ to $K$ such that
\[
G(x) = (g_1(x), \ldots, g_m(x))
\]
for all $x \in U$. We call $G$ an \emph{analytic mapping} if the $m$
functions $g_1, \ldots, g_m$ are analytic in $U$.

\begin{theorem*}[2.1.5 (Composition Theorem)]
Let $U$ and $U'$ be open subsets of $K^n$ and $K^m$ respectively, let
$G: U \to U'$ be an analytic mapping, and let $f: U' \to K$ be an
analytic function. Then the composite $f \circ G$ is an analytic function
in $U$.
\end{theorem*}

\textit{References for proof.} For the real case, see the last paragraph
of Sec.~2, Ch.~2 of [BMA48]. For the complex case, see Theorem 5 of
Sec.~A, Ch.~I of [GRO65]. $\square$

If $F$ is an analytic mapping from $U \subseteq K^n$ into $K^m$,
$F(x) = (f_1(x), \ldots, f_m(x))$, and if $p$ is a point of $U$, then we
denote the Jacobian matrix $\left(\dfrac{\partial f_i}{\partial x_j}(p)\right)$
of $F$ at $p$ by $J_F(p)$. We can now state a couple of fundamental
theorems on the local properties of analytic mappings. These theorems,
the inverse mapping and the implicit mapping theorems, generalize the
familiar inverse function and implicit function theorems from the
elementary calculus.

\begin{theorem*}[2.1.6 (Inverse Mapping Theorem)]
Let $F$ be an analytic mapping from the open subset $U$ of $K^n$ into
$K^n$, and let $p$ be a point of $U$. Suppose that $J_F(p)$ is
invertible. Then there is a neighborhood $V \subseteq U$ of $p$ in which
$F$ is invertible. That is, the set $V' := F(V)$ is open in $K^n$, and
there is an analytic map $G: V' \to V$ such that, for all $x$ in $V$ and
$y$ in $V'$,
\[
y = F(x) \text{ if and only if } x = G(y).
\]
\end{theorem*}

\textit{References.} The real case of the theorem is stated without
proof in the appendix on Calculus in [HIR76]. (Note that the linear map
$Df_p : \R^n \to \R^n$, which is mentioned in the statement of the
theorem in [HIR76], has matrix $J_F(p)$ with respect to the standard
basis of unit coordinate vectors of $\R^n$. Hence, the mapping $Df_p$ is
invertible if and only if the matrix $J_F(p)$ is invertible.)

The complex version of the theorem appears in [GRO65]: it is Theorem 7
of Sec.~B, Ch.~I of this book. $\square$

\begin{theorem*}[2.1.7 (Implicit Mapping Theorem)]
Let $0 \le s < n$, let $U$ be an open subset of $K^n$, and let
$F: U \to K^{n-s}$ be an analytic map with component functions
$f_{s+1}, \ldots, f_n$. Let $p = (p_1, \ldots, p_n)$ be a point of $U$
and suppose $F(p) = 0$. Suppose that the $(n-s) \times (n-s)$ square
submatrix of $J_F(p)$ consisting of the last $n-s$ rows and columns of
$J_F(p)$ is invertible. Then there are neighborhoods $V \subseteq K^s$
of $(p_1, \ldots, p_s)$ and $W \subseteq K^{n-s}$ of
$(p_{s+1}, \ldots, p_n)$, with $V \times W \subseteq U$, and analytic
functions $\psi_{s+1}, \ldots, \psi_n$ from $V$ into $K$, such that for
all $(x_1, \ldots, x_s)$ in $V$ and all $(x_{s+1}, \ldots, x_n)$ in $W$,
\[
F(x_1, \ldots, x_n) = 0 \text{ iff }
x_{s+1} = \psi_{s+1}(x_1, \ldots, x_s),\ \ldots,\
x_n = \psi_n(x_1, \ldots, x_s).
\]
\end{theorem*}

This result can be proved using the inverse mapping theorem
(cf.~[BUC56]).

Up to now, the theories of real and complex analytic functions have been
developed in parallel. We now mention a couple of differences between
the two theories. The following theorem holds for complex but not real
analytic functions:

\begin{theorem*}[2.1.8]
Let $U$ be an open subset of $\C^n$ and let $f$ be a complex-valued
function defined in $U$. Suppose that $f$ is continuous in $U$ and that
each partial derivative $\dfrac{\partial f}{\partial z_i}$ exists and is
continuous in $U$. Then $f$ is analytic in $U$.
\end{theorem*}

\textit{References for Proof.} The result follows from the $n$-variable
analogue of Theorem 53 in Sec.~9-3 of [KAP66] (note that the definition
of holomorphicity in [KAP66] is equivalent to the conditions in the
hypothesis of Theorem 2.1.8).

We include one more theorem that is valid for holomorphic functions but
which does not have a real analogue.

\begin{theorem*}[2.1.9]
Let $U$ be an open subset of $\C^n$ and let $f$ be holomorphic in $U$.
Let $c$ be a point of $U$, let $r = (r_1, \ldots, r_n)$, where each
$r_i > 0$, and suppose that $\Delta(c;r) \subseteq U$. Then the power
series expansion of $f$ about $c$ is absolutely convergent in
$\Delta(c;r)$.
\end{theorem*}

\textit{Reference for proof.} See Theorem 3 in Ch.~II of [BMA48].

\begin{remark}
In a real analogue of the above theorem one would presumably have boxes
$B(c;r)$ in place of polydiscs $\Delta(c;r)$. That there is no real
analogue of the theorem, however, is shown by the following example: the
function $f(x) = \dfrac{1}{1 + x^2}$ is analytic on the whole real line,
but the interval of convergence of the power series expansion for $f$
about $0$
\[
1 - x^2 + x^4 - x^6 + \cdots
\]
is $|x| < 1$.
\end{remark}

By a \emph{zero} of a holomorphic function $f(z_1, \ldots, z_n)$ is meant
a point $p = (p_1, \ldots, p_n)$ such that $f(p) = 0$. It is noted in
Sec.~9-7 of [KAP66] that if $n \ge 2$, then a holomorphic function
$f(z_1, \ldots, z_n)$ can have no isolated zeros. Further information on
the set of zeros is given by the Weierstrass preparation theorem. We
need a definition before we can state the theorem. Let $z$ denote the
$(n-1)$-tuple $(z_1, \ldots, z_{n-1})$.

\begin{definition}
Let
\begin{equation}
h(z, z_n) = a_0(z) z_n^m + a_1(z) z_n^{m-1} + \cdots + a_m(z)
\tag{2.1.6}
\end{equation}
where each $a_i(z)$ is holomorphic in the polydisc
$\Delta_1 \subseteq \C^{n-1}$ about $0$. Then $h(z, z_n)$ is called a
\emph{pseudopolynomial} (in $\Delta_1$). If $a_0$ is not identically
zero, then the \emph{degree} of $h(z, z_n)$ is $m$. If $a_0 \equiv 1$ and
$a_i(0) = 0$ for each $i$, $1 \le i \le m$, then $h(z, z_n)$ is called a
\emph{Weierstrass polynomial} (in $\Delta_1$).
\end{definition}

\begin{theorem*}[2.1.10 (Weierstrass preparation theorem)]
Let $f(z, z_n)$ be holomorphic in the polydisc
$\Delta_1 \times \Delta(0;\epsilon) \subseteq \C^{n-1} \times \C$, let
$z_n = 0$ be a root of $f(0, z_n)$ of multiplicity $m \ge 1$, and assume
that $f(0, z_n) \ne 0$ for $0 < |z_n| \le \epsilon$. Then there is a
polydisc $\Delta_2 \subseteq \Delta_1$ about $0$ in $\C^{n-1}$, a
function $u(z, z_n)$ holomorphic and non-vanishing in
$\Delta' := \Delta_2 \times \Delta(0;\epsilon)$, and a Weierstrass
polynomial $h(z, z_n)$ in $\Delta_2$, of the form~(2.1.6) (with
$a_0 \equiv 1$), such that
\begin{equation}
f(z, z_n) = u(z, z_n) h(z, z_n)
\tag{2.1.7}
\end{equation}
for all $(z, z_n) \in \Delta'$, and such that for each fixed
$z \in \Delta_2$, all the $m$ roots of $h(z, z_n)$ are contained in the
disc $\Delta(0;\epsilon)$.
\end{theorem*}

\textit{References for proof.} We refer the reader to the proof of the
Weierstrass preparation theorem (Theorem 62) in Ch.~9 of [KAP66].
(Although there are minor respects in which our statement of the
Weierstrass preparation theorem differs from that of Kaplan, the careful
reader will note that our theorem follows from the proof of Theorem 62
of [KAP66].)

\section{Submanifolds of Euclidean Space}

In this section we develop the concept of a submanifold of real
$n$-space $\R^n$.

We first give a summary of the material presented in this section. An
\emph{$s$-submanifold} of $\R^n$ is a set $S$ which ``looks locally like
Euclidean $s$-space $\R^s$''; that is, for every point $p$ of $S$, there
is an (analytic) coordinate system about $p$ with respect to which $S$
is locally the intersection of some $n-s$ coordinate hyperplanes. For
example, the $(n-1)$-sphere
$S^{n-1} = \{ x = (x_1, \ldots, x_n) \in \R^n : \sum x_i^2 = 1 \}$ is an
$(n-1)$-submanifold of $\R^n$. The first part of this section gives a
rigorous definition of submanifold, and shows that it corresponds to the
above intuitive notion.

Nonempty open subsets of $\R^n$ are submanifolds~--- in fact, the
nonempty open subsets are the $n$-dimensional submanifolds. But
lower-dimensional submanifolds are not open. One can, however, define in
a natural way the notion of an analytic function from an $s$-submanifold
$S$ into $\R$, where $0 \le s \le n$. A function $f: S \to \R$ is said to
be analytic if for every point $p$ of $S$, there is a coordinate system
about $p$ with respect to which $S$ looks locally like $\R^s$, and with
respect to which $f$ looks locally like an analytic function from $\R^s$
to $\R$. The second part of this section gives a rigorous definition of
an analytic function on a submanifold, and proves that the definition is
independent of any particular coordinate system for the submanifold.

The third and last part of this section gives a couple of results on the
creation of new submanifolds from old. In particular, it is shown that
the graph of an analytic function defined on a submanifold of $\R^n$ is
a submanifold of $\R^{n+1}$.

We now begin to look at the material described above in detail. Before
giving a formal definition of submanifold we define the notion of a
regular point of an analytic mapping.

\begin{definition}
Let $U \subseteq \R^n$ be open and let $F: U \to \R^m$, $m \le n$, be an
analytic map. For $x = (x_1, \ldots, x_n) \in U$ let
$F(x) = (F_1(x), \ldots, F_m(x))$. The point $p$ of $U$ is said to be a
\emph{regular point} of $F$ if the rank of the Jacobian matrix
$J_F(p) = \left(\dfrac{\partial F_i}{\partial x_j}(p)\right)$ of $F$ at
$p$ is equal to $m$.
\end{definition}

\begin{example}
Let $F: \R^3 \to \R$ be defined by $F(x, y, z) = x^2 + y^2 + z^2 - 1$.
Then $J_F = (2x, 2y, 2z)$, so every point of $\R^3$ other than the
origin is a regular point of $F$.
\end{example}

\begin{definition}
The nonempty subset $S$ of $\R^n$ is an \emph{analytic $s$-dimensional
submanifold of $\R^n$} (or \emph{$C^\omega$ $s$-submanifold}, or
\emph{$C^\omega$ smooth}, for short), where $0 \le s \le n$, if for
each point $p$ of $S$ there is a neighborhood $W \subseteq \R^n$ of $p$
and an analytic map $F: W \to \R^{n-s}$ which has $p$ as a regular
point, such that
\[
S \cap W = \{ x \in W : F(x) = 0 \}.
\]
\end{definition}

We remark that there is a notion of an $s$-submanifold of $\R^n$ of
class $C^r$, defined as above in terms of maps $F$ which are required
to be of class $C^r$ (i.e.\ to possess continuous partial derivatives
through the order $r$). Here, $r$ is allowed to be any non-negative
integer, $\infty$, or $\omega$ (meaning analytic). The only kind of
submanifold we shall consider is the analytic kind. Thus we shall
henceforth omit the term `analytic' when referring to submanifolds: all
submanifolds will be understood to be analytic.

\begin{example}
Let $S^2 = \{ (x, y, z) \in \R^3 : x^2 + y^2 + z^2 = 1 \}$ be the unit
sphere in $\R^3$. For each point $p$ in $S^2$ we may take $W = \R^3$ and
$F: W \to \R$ to be the map $F(x, y, z) = x^2 + y^2 + z^2 - 1$. As
noted above, $F$ is regular at every point $p \ne 0$, and hence at
every point $p$ of $S^2$. Thus $S^2$ is a 2-submanifold of $\R^3$.
\end{example}

\begin{definition}
Let $U$ and $V$ be open subsets of $\R^n$. A homeomorphism
$\Phi: U \to V$ such that both $\Phi$ and $\Phi^{-1}$ are analytic maps
is called an \emph{analytic isomorphism}.
\end{definition}

\begin{definition}
Let $U$ and $V$ be open subsets of $\R^n$ and let $\Phi: U \to V$ be an
analytic isomorphism. Then $\Phi$ is called a \emph{coordinate system in
$U$}. Suppose that $p \in U$, $0 \in V$, and $\Phi(p) = 0$. Then $\Phi$
is called a \emph{coordinate system (in $U$) about $p$}.
\end{definition}

The next theorem expresses the intuitive idea that an $s$-submanifold of
$\R^n$, $0 \le s \le n$, is a set which ``looks locally like Euclidean
$s$-space''.

\begin{theorem*}[2.2.1]
The nonempty subset $S$ of $\R^n$ is an $s$-submanifold of $\R^n$,
$0 \le s \le n$, if and only if for every point $p$ of $S$ there is a
neighborhood $U \subseteq \R^n$ of $p$ and a coordinate system
$\Phi: U \to V$, $\Phi = (\phi_1, \ldots, \phi_n)$, about $p$ such that
\begin{equation}
S \cap U = \{ x \in U : \phi_{s+1}(x) = 0, \ldots, \phi_n(x) = 0 \}.
\tag{2.2.1}
\end{equation}
\end{theorem*}

\begin{remarks}
Let $0 \le s \le n$ and let $T$ be the $s$-dimensional linear subspace
\begin{equation}
T = \{ y = (y_1, \ldots, y_n) \in \R^n : y_{s+1} = 0, \ldots, y_n = 0 \}
\tag{2.2.2}
\end{equation}
of $\R^n$. Then $\R^s$ may be identified with $T$ under the natural
identification mapping
$\iota(y_1, \ldots, y_s) = (y_1, \ldots, y_s, 0, \ldots, 0)$ (we write
$\R^s \equiv T$). By Theorem 2.2.1, the nonempty subset $S$ of $\R^n$ is
an $s$-submanifold of $\R^n$ if and only if $S$ ``looks locally like the
subspace $T \equiv \R^s$ of $\R^n$'' (i.e.\ for each point $p$ of $S$
there is a neighborhood $U \subseteq \R^n$ of $p$ and a coordinate system
$\Phi: U \to V$ about $p$ such that $\Phi(S \cap U) = T \cap V$.)
\end{remarks}

\begin{proof}[Proof of 2.2.1]
Suppose that $S$ is an $s$-submanifold of $\R^n$. Let $p$ be a point of
$S$. Then there is a nbd $W \subseteq \R^n$ of $p$ and an analytic map
$F: W \to \R^{n-s}$ having $p$ as a regular point such that
$S \cap W = \{ x \in W : F(x) = 0 \}$. If $s = n$ then set $U = W$ and
$\Phi(x) = x - p$: note that~(2.2.1) holds. Assume that $s < n$. Write
$F(x) = (f_{s+1}(x), \ldots, f_n(x))$. As the Jacobian matrix $J_F(p)$
of $F$ at $p$ has rank $n-s$, we may assume after a renumbering of the
coordinates that the $(n-s) \times (n-s)$ submatrix
\[
\left( \frac{\partial f_i}{\partial x_j}(p) \right)_{\substack{s+1 \le i \le n \\ s+1 \le j \le n}}
\]
of $J_F(p)$ is invertible. Define $\Phi: W \to \R^n$,
$\Phi = (\phi_1, \ldots, \phi_n)$, as follows:
\[
\phi_i(x) =
\begin{cases}
x_i - p_i, & 1 \le i \le s \\
f_i(x), & s+1 \le i \le n.
\end{cases}
\]
Then $\Phi$ is an analytic map, and the Jacobian matrix $J_\Phi(p)$ of
$\Phi$ at $p$ is invertible. Hence, by the inverse mapping theorem
(Theorem 2.1.6), there is a neighborhood $U \subseteq W$ of $p$ in which
$\Phi$ is invertible (i.e.\ $V := \Phi(U)$ is open in $\R^n$, and the
restriction of $\Phi$ to $U$ has an analytic inverse which maps $V$ onto
$U$). By definition of $\Phi$, (2.2.1) holds. The other direction of the
theorem is obvious.
\end{proof}

Now nonempty open subsets of $\R^n$ are submanifolds~--- in fact the
nonempty open subsets are the $n$-dimensional submanifolds. But
lower-dimensional submanifolds are not open. One can, however, define in
a natural way the notion of an analytic function from an $s$-submanifold
$S$ into $\R$, where $0 \le s \le n$. Having the notion of chart at
one's disposal facilitates this definition. Let $S$ be a subset of
$\R^n$. A \emph{chart for $S$} is a homeomorphism from an open subset of
$S$ onto an open subset of $\R^s$, for some $s$, $0 \le s \le n$. Let
$S$ now be an $s$-dimensional submanifold, $0 \le s \le n$, and let $p$
be a point of $S$. Let $U \subseteq \R^n$ be a neighborhood of $p$, and
let $\Phi: U \to V$, $\Phi = (\phi_1, \ldots, \phi_n)$, be a coordinate
system about $p$ such that $\Phi(S \cap U) = T \cap V$, where $T$ is
given by~(2.2.2). As remarked following the statement of Theorem 2.2.1,
$T \equiv \R^s$ under the identification mapping $\iota: \R^s \to T$:
let $T \cap V$ correspond to the neighborhood $W$ of $0$ in $\R^s$ under
this identification (written $T \cap V \equiv W$). The mapping
$\phi: S \cap U \to W$ given by
\[
\phi(x) = (\phi_1(x), \ldots, \phi_s(x))
\]
is a homeomorphism with analytic inverse $\phi^{-1}$\footnote{$\phi^{-1}$
is the composite of $\Phi^{-1}$ and the restriction of $\iota$ to $W$.},
and is called the \emph{chart for $S$ corresponding to $\Phi$}.

\begin{definition}
Let $S$ be an $s$-submanifold of $\R^n$, $0 \le s \le n$, let $T$ be
given by~(2.2.2), and let $f: S \to \R$ be a function. Then $f$ is said
to be \emph{analytic (in $S$)} if for each point $q$ of $S$ there is a
neighborhood $U$ of $q$ and a coordinate system $\Phi: U \to V$ about
some point $p$ of $S \cap U$, such that
\begin{enumerate}
\item[(a)] $\Phi(S \cap U) = T \cap V$; and
\item[(b)] $f \circ \phi^{-1}: W \to \R$ is analytic, where $W$ is the
neighborhood of $0$ in $\R^s$ for which $T \cap V \equiv W$, and
$\phi: S \cap U \to W$ is the chart for $S$ corresponding to $\Phi$.
\end{enumerate}
\end{definition}

Condition (b) in the above definition can be paraphrased, ``$f$ is
analytic with respect to the coordinate system $\Phi$''. The following
theorem states that if $f: S \to \R$ is analytic then, for each point
$p$ of $S$, $f$ is analytic with respect to any coordinate system
$\Phi$ about $p$ satisfying condition (a) of the above definition.

\begin{theorem*}[2.2.2]
Let $S$ be an $s$-submanifold of $\R^n$, $0 \le s \le n$, let $T$ be
given by~(2.2.2), and let $f: S \to \R$ be an analytic function. Let $p$
be a point of $S$ and let $\Phi: U \to V$ be a coordinate system about
$p$ such that $\Phi(S \cap U) = T \cap V \equiv W$. Then
$f \circ \phi^{-1}: W \to \R$ is analytic, where $\phi: S \cap U \to W$
is the chart for $S$ corresponding to $\Phi$.
\end{theorem*}

\begin{proof}
Let $w \in W$. Then $w = \phi(q)$ for some point $q$ of $S \cap U$. By
definition there is a neighborhood $U'$ of $q$ and a coordinate system
$\Psi: U' \to V'$ such that $\Psi(S \cap U') = T \cap V' \equiv W'$ and
$f \circ \psi^{-1}: W' \to \R$ is analytic, where $\psi: S \cap U' \to W'$
is the chart for $S$ corresponding to $\Psi$. Let $\pi: \R^n \to \R^s$
be the projection $\pi(y_1, \ldots, y_n) = (y_1, \ldots, y_s)$. Then we
have
$f \circ \phi^{-1}(y_1, \ldots, y_s)
= f \circ \psi^{-1}\bigl((\pi \circ \Psi \circ \phi^{-1})(y_1, \ldots, y_s)\bigr)$
for every point $(y_1, \ldots, y_s)$ of the neighborhood
$\phi(S \cap U \cap U')$ of $w$. Therefore, by the composition theorem
(Theorem 2.1.5), $f \circ \phi^{-1}$ is analytic near $w$. Hence
$f \circ \phi^{-1}$ is analytic in $W$.
\end{proof}

Finally, a couple of results on the creation of new submanifolds from
old.

\begin{theorem*}[2.2.3]
Let $S$ be an $s$-submanifold of $\R^{n-1}$, $0 \le s \le n-1$, and let
$f: S \to \R$ be an analytic function. Then the graph of $f$ is an
$s$-submanifold of $\R^n$.
\end{theorem*}

\begin{proof}
Let $G$ be the graph of $f$ and let $(p, p_n)$ be a point of $G$, where
$p = (p_1, \ldots, p_{n-1})$. We shall find a coordinate system about
$(p, p_n)$ in $\R^n$ with respect to which $G$ is locally the
intersection of the last $n-s$ coordinate hyperplanes. By Theorem 2.2.1
there is a neighborhood $U \subseteq \R^{n-1}$ of $p$ and a coordinate
system $\Phi: U \to V$ about $p$ such that $\Phi(S \cap U) = T \cap V$,
where $T$ is given by~(2.2.2). Let $W$ be the neighborhood of $0$ in
$\R^s$ such that $W \equiv T \cap V$ and let $\phi: S \cap U \to W$ be
the chart for $S$ corresponding to $\Phi$. Let $\pi: \R^{n-1} \to \R^s$
be the projection $\pi(y_1, \ldots, y_{n-1}) = (y_1, \ldots, y_s)$. Then
$\pi \circ \Phi$ is an analytic map from $U$ into $\R^s$. Hence, as $W$
is a neighborhood of $0$ in $\R^s$, the set
$U' := (\pi \circ \Phi)^{-1}(W)$ is a neighborhood of $p$ in $\R^{n-1}$
with $U' \subseteq U$. Let $V' = \Phi(U')$ and define a map
$\Psi: U' \times \R \to V' \times \R$ by
\[
\Psi(x, x_n) = \bigl( \Phi(x),\
x_n - f \circ \phi^{-1}((\pi \circ \Phi)(x)) \bigr),
\]
for $(x, x_n) \in U' \times \R$. Then $\Psi$ is an analytic map (the
$n$-th component of $\Psi$ is analytic because $f \circ \phi^{-1}$ is
analytic in $W$ by Theorem 2.2.2). In fact, one can verify that $\Psi$
is invertible, with analytic inverse $\Psi^{-1}$ given by
\[
\Psi^{-1}(y, y_n) = \bigl( \Phi^{-1}(y),\
y_n + f \circ \phi^{-1}(y_1, \ldots, y_s) \bigr),
\]
for $y = (y_1, \ldots, y_{n-1}) \in V'$ and $y_n \in \R$. Thus $\Psi$ is
a coordinate system about $(p, p_n)$. Where $\psi_1, \ldots, \psi_n$ are
the components of $\Psi$, we have
\[
G \cap (U' \times \R) = \{ (x, x_n) \in U' \times \R :
\psi_{s+1}(x, x_n) = 0, \ldots, \psi_n(x, x_n) = 0 \}.
\]
By Theorem 2.2.1, $G$ is an $s$-submanifold of $\R^n$.
\end{proof}

\begin{theorem*}[2.2.4]
Let $S$ be an $s$-submanifold of $\R^{n-1}$, $0 \le s \le n-1$, and let
$f$ and $g$ be continuous functions from $S$ into $\R$ (also allowed are
$f \equiv -\infty$ or $g \equiv +\infty$) with $f < g$. Let
$R = \{ (x, x_n) \in S \times \R : f(x) < x_n < g(x) \}$. Then $R$ is an
$(s+1)$-submanifold of $\R^n$.
\end{theorem*}

\begin{proof}
Let $(p, p_n)$ be a point of $R$. By Theorem 2.2.1 there is a
neighborhood $U \subseteq \R^{n-1}$ of $p$ and a coordinate system
$\Phi: U \to V$ about $p$, $\Phi = (\phi_1, \ldots, \phi_{n-1})$, such
that
\[
S \cap U = \{ x \in U : \phi_{s+1}(x) = 0, \ldots, \phi_{n-1}(x) = 0 \}.
\]
Choose $\epsilon > 0$ such that $f(p) < p_n - \epsilon < p_n + \epsilon < g(p)$.
By continuity of $f$ and $g$, there exists a neighborhood
$U' \subseteq U$ of $p$ in $\R^{n-1}$ such that $f(x) < x_n < g(x)$ for
every $x \in S \cap U'$ and every $x_n \in (p_n - \epsilon, p_n + \epsilon)$.
Let $I = (p_n - \epsilon, p_n + \epsilon)$, let
$J = (-\epsilon, +\epsilon)$, and let $V' = \Phi(U')$. Define
$\Psi: U' \times I \to V' \times J$ by
$\Psi(x, x_n) = (\Phi(x), x_n - p_n)$. Then $\Psi$ is an analytic map,
with analytic inverse $\Psi^{-1}$ given by
$\Psi^{-1}(y, y_n) = (\Phi^{-1}(y), y_n + p_n)$. Where
$\psi_1, \ldots, \psi_n$ are the components of $\Psi$, we have
\[
R \cap (U' \times I) = \{ (x, x_n) \in U' \times I :
\psi_{s+1}(x, x_n) = 0, \ldots, \psi_{n-1}(x, x_n) = 0 \}.
\]
Even though $\psi_n(x, x_n) = 0$ is not included amongst the local
defining equations for $R$, by an obvious analogue of Theorem 2.2.1,
$R$ is an $(s+1)$-submanifold of $\R^n$.
\end{proof}

\section{Miscellaneous Results}

\begin{definition}
Let $K = \R$ or $\C$. Let $U$ be an open subset of $K^n$ and let
$f: U \to K$ be an analytic function. Let $p$ be a point of $U$. If some
partial derivative of $f$ of non-negative order does not vanish at $p$
then we say that $f$ has \emph{order $k$ at $p$}, and write
$\ord_p f = k$, provided that $k$ is the least non-negative integer such
that some partial derivative of $f$ of order $k$ does not vanish at $p$;
we also say that $f$ has \emph{order $k_i$ at $p$ in $x_i$}, provided
that $k_i$ is the least non-negative integer such that some partial
derivative of $f$ of order $k_i$ in $x_i$ does not vanish at $p$. If,
on the other hand, all partial derivatives of $f$ of all orders vanish
at $p$, then we say that $f$ has \emph{order $\infty$ at $p$}, and write
$\ord_p f = \infty$; we also say that $f$ has order $\infty$ at $p$ in
each $x_i$.
\end{definition}

\begin{theorem*}[2.3.1]
Let $K = \R$ or $\C$. Let $U \subseteq K^n$ and $V \subseteq K^m$ be
open sets, let $G: U \to V$ be an analytic mapping, and let
$f: V \to K$ be an analytic function. Then, for every point $p$ of $U$,
\[
\ord_{G(p)} f \le \ord_p f \circ G.
\]
\end{theorem*}

\begin{proof}
Let the coordinates of $K^n$ be $(x_1, \ldots, x_n)$ and let those of
$K^m$ be $(y_1, \ldots, y_m)$. Let $h = f \circ G$. We prove by induction
on $k$ that every partial derivative of $h$ of order $k$ is a finite sum
of terms of the form
\[
(P \circ G) Q
\]
where $P = P(y_1, \ldots, y_m)$ is a partial derivative of
$f = f(y_1, \ldots, y_m)$ of order $\le k$ and $Q$ is an analytic
function. This is true for $k = 0$ by definition of $h$. Assume that the
above proposition is true for $k \ge 0$. Let
$R = \dfrac{\partial^{k+1} h}{\partial x_1^{i_1} \cdots \partial x_n^{i_n}}$
be a partial derivative of $h$ of order $k+1$. As some $i_s > 0$, we
have $R = \dfrac{\partial S}{\partial x_s}$, where $S$ is a partial
derivative of $h$ of order $k$. By the induction hypothesis, $S$ is a
finite sum of terms of the form $(P \circ G) Q$, where $P$ is a partial
derivative of $f$ of order $\le k$. (Let $g_1, \ldots, g_m$ be the
component functions of $G$.) As
\[
\frac{\partial (P \circ G) Q}{\partial x_s}
= \frac{\partial (P \circ G)}{\partial x_s} \cdot Q
+ (P \circ G) \frac{\partial Q}{\partial x_s}
\]
and
\[
\frac{\partial (P \circ G)}{\partial x_s}
= \left( \frac{\partial P}{\partial y_1} \circ G \right)
\frac{\partial g_1}{\partial x_s} + \cdots
+ \left( \frac{\partial P}{\partial y_m} \circ G \right)
\frac{\partial g_m}{\partial x_s},
\]
it follows that $\dfrac{\partial (P \circ G) Q}{\partial x_s}$ is a
finite sum of terms of the form $(P' \circ G) Q'$, where $P'$ is a
partial derivative of $f$ of order $\le k + 1$. Hence $R$ is such a
sum. We have shown that the required proposition holds for $k + 1$.
The proof by induction is complete.

Let $p$ be a point of $U$ and let the order of $f$ at $G(p)$ be $l$. If
$l = 0$ then clearly $\ord_p h \ge 0$, so assume $l > 0$. Let $R$ be a
partial derivative of $h$ of order $k < l$. Then $R$ is a finite sum of
terms of the form $(P \circ G) Q$, where $P$ is a partial derivative of
$f$ of order $\le k$. But $P(G(p)) = 0$ for every partial derivative
$P$ of $f$ of order $\le k$, as $\ord_{G(p)} f = l > k$. Hence
$R(p) = 0$. It follows that $\ord_p h \ge l$.
\end{proof}

\begin{theorem*}[2.3.2 (Root continuity principle)]
Let $f(z) = f_0 z^d + f_1 z^{d-1} + \cdots + f_d$ be a polynomial in
$\C[z]$, with $f_0 = f_1 = \cdots = f_{l-1} = 0$ and $f_l \ne 0$, for
some $l$, $0 \le l \le d$. Let $\alpha$ be a root of $f(z)$ of
multiplicity $m$ and let $C$ be a circle in the complex plane centered
at $\alpha$, of radius $\epsilon > 0$, such that $f \ne 0$ in the
punctured disc $0 < |z - \alpha| \le \epsilon$. Then there is a number
$\delta > 0$ such that if $g_0, g_1, \ldots, g_d \in \C$ and
$|g_j - f_j| < \delta$ for $0 \le j \le d$, then the polynomial
\begin{equation}
g(z) = g_0 z^d + g_1 z^{d-1} + \cdots + g_d
\tag{2.3.1}
\end{equation}
has exactly $m$ roots inside $C$.
\end{theorem*}

\begin{proof}
There exists $\delta > 0$ such that if $g_0, \ldots, g_d \in \C$,
$|g_j - f_j| < \delta$ for $0 \le j \le d$, and $g(z)$ is given
by~(2.3.1), then $|g(z) - f(z)| < |f(z)|$ for $|z| = \epsilon$. Let
$g_0, \ldots, g_d \in \C$ and let $|g_j - f_j| < \delta$ for
$0 \le j \le d$. Let $g(z)$ be defined by~(2.3.1) and let
$f_s(z) = f(z) + s(g(z) - f(z))$, for $0 \le s \le 1$. Now if
$0 \le s \le 1$ then
\[
\begin{aligned}
|f_s(z)| &= |f(z) + s(g(z) - f(z))| \ge \bigl| |f(z)| - |s(g(z) - f(z))| \bigr| \\
&= \bigl| |f(z)| - s|g(z) - f(z)| \bigr| \\
&> 0
\end{aligned}
\]
for $|z| = \epsilon$. Hence, if $0 \le s \le 1$, then $f_s(z)$ has a
total number $m_s$ of zeros inside $C$ given by
\begin{equation}
m_s = \frac{1}{2\pi i} \int_C \frac{f'_s(\zeta) \, d\zeta}{f_s(\zeta)}.
\tag{2.3.2}
\end{equation}
But the integral on the right-hand side of~(2.3.2) is continuous in $s$
and is integer-valued, and hence is constant. Therefore $m_0 = m_1$.
\end{proof}

Let $R$ be an integral domain and let
\[
\begin{aligned}
A(x) &= a_0 x^m + a_1 x^{m-1} + \cdots + a_m, \\
B(x) &= b_0 x^n + b_1 x^{n-1} + \cdots + b_n
\end{aligned}
\]
be polynomials in $R[x]$, not both constant, of degrees $m$ and $n$
respectively. The \emph{Sylvester matrix} of $A$ and $B$ is the
$m+n$ by $m+n$ matrix
\[
M =
\begin{pmatrix}
a_0 & a_1 & \cdot & \cdot & \cdot & a_m & & & \\
& a_0 & a_1 & \cdot & \cdot & \cdot & a_m & & \\
& & \cdot & \cdot & \cdot & \cdot & \cdot & \cdot & \\
& & & a_0 & a_1 & \cdot & \cdot & \cdot & a_m \\
b_0 & b_1 & \cdot & \cdot & \cdot & b_n & & & \\
& b_0 & b_1 & \cdot & \cdot & \cdot & b_n & & \\
& & \cdot & \cdot & \cdot & \cdot & \cdot & \cdot & \\
& & & b_0 & b_1 & \cdot & \cdot & \cdot & b_n
\end{pmatrix}
\]
The \emph{resultant} of $A$ and $B$, $\res(A, B)$, is defined by
\[
\res(A, B) = \det(M).
\]
Suppose now that the characteristic of $R$ is zero and that $A$ is not
constant. Then $m a_0 \ne 0$, and so the degree of the derivative
\[
A'(x) = m a_0 x^{m-1} + (m-1) a_1 x^{m-2} + \cdots + a_{m-1}
\]
of $A(x)$ is $m - 1$. Consider the Sylvester matrix of $A$ and $A'$.
The first column has two nonzero entries, $a_0$ and $m a_0$. Hence we
can factor out $a_0$ from $\res(A, A')$. The \emph{discriminant} of
$A$, $\discr(A)$, is defined by the equation
\[
a_0 \discr(A) = (-1)^{m(m-1)/2} \res(A, A').
\]
Let $K$ be the quotient field of $R$ and let $\Omega$ be an algebraic
closure of $K$. Then we can completely factor $A(x)$ and $B(x)$ into
linear factors over $\Omega$ thus:
\begin{align}
A(x) &= a_0 (x - \alpha_1) \cdots (x - \alpha_m) \tag{2.3.3} \\
B(x) &= b_0 (x - \beta_1) \cdots (x - \beta_n). \tag{2.3.4}
\end{align}
The following expressions for $\res(A, B)$ and $\discr(A)$ are
well-known:
\begin{align}
\res(A, B) &= a_0^n b_0^m \prod_{i, j} (\alpha_i - \beta_j)
\tag{2.3.5} \\
\discr(A) &= a_0^{2m-2} \prod_{i < j} (\alpha_i - \alpha_j)^2.
\tag{2.3.6}
\end{align}

\begin{theorem*}[2.3.3]
Let $A(x)$ and $B(x)$ be non-constant polynomials over the integral
domain $R$ of characteristic zero. Let $C(x) = A(x) B(x)$. Then
\[
\discr(C) = \discr(A) (\res(A, B))^2 \discr(B).
\]
\end{theorem*}

\begin{proof}
Let $K$ be the quotient field of $R$ and let $\Omega$ be an algebraic
closure of $K$. Factor $A(x)$ and $B(x)$ completely over $\Omega$ as
in~(2.3.3) and~(2.3.4). Let
$\gamma_1 = \alpha_1, \ldots, \gamma_m = \alpha_m,\
\gamma_{m+1} = \beta_1, \ldots, \gamma_{m+n} = \beta_n$. Let
$c_0 = a_0 b_0$ and $l = m + n$. Then
\[
C(x) = c_0 (x - \gamma_1) \cdots (x - \gamma_l).
\]
Therefore
\[
\begin{aligned}
\discr(C) &= c_0^{2l - 2} \prod_{1 \le i < j \le l} (\gamma_i - \gamma_j)^2 \\
&= \left( a_0^{2m-2} \prod_{1 \le i < j \le m} (\alpha_i - \alpha_j)^2 \right)
\cdot
\left( a_0^{2n} b_0^{2m}
\prod_{\substack{1 \le i \le m \\ 1 \le j \le n}} (\alpha_i - \beta_j)^2 \right) \\
&\quad \cdot \left( b_0^{2n-2} \prod_{1 \le i < j \le n} (\beta_i - \beta_j)^2 \right) \\
&= \discr(A) (\res(A, B))^2 \discr(B).
\end{aligned}
\]
\end{proof}

The \emph{sum-norm} (or just \emph{norm} for short) of an integral
polynomial $A$ in $r$ variables is the sum of the absolute values of the
integer coefficients of $A$. The \emph{max-norm} of such a polynomial is
the maximum of the absolute values of the integer coefficients. We
denote the sum-norm by $|A|_1$, and the max-norm by $|A|_\infty$. The
following theorem leads to a bound for the coefficients of the factors
of an $r$-variate integral polynomial.

\begin{theorem*}[2.3.4 (Gelfond)]
Let $A_1, \ldots, A_k$ be $r$-variate polynomials with complex
coefficients, and let $A = A_1 \cdots A_k$. Let $n_i$ be the degree of
$A$ in the $i$-th variable, assume $n_i > 0$ for all $i$, and let
$m = \sum_{i=1}^{r} n_i$. Then
\[
\prod_{j=1}^{k} |A_j|_\infty \le 2^{m-r/2} |A|_1.
\]
\end{theorem*}

\begin{proof}
See [GEL60], pp 135--139.
\end{proof}

\begin{theorem*}[Corollary 2.3.5]
Let $A$ be a non-zero integral polynomial in $r$ variables, and let $B$
be a factor of $A$. Let $n$ be a bound for the degree of $A$ in each
variable. Then
\[
|B|_1 \le (n+1)^r \, 2^{rn - r/2} |A|_1.
\]
\end{theorem*}

\begin{proof}
As $B$ has at most $(n+1)^r$ terms, we have
\[
|B|_1 \le (n+1)^r |B|_\infty.
\]
The corollary now follows from the theorem.
\end{proof}

\section*{Chapter Three: Reduced Projection Map for Cylindrical Algebraic Decomposition}

In this chapter we launch the main thrust of the thesis. Section 1
contains a review of the essential aspects of the cylindrical algebraic
decomposition (cad) algorithm. The main focus of the review is the
projection operation used in the cad algorithm. In Section 2 we
introduce a new projection map, which is a reduced version of the
original. Theorems are presented which suggest the usefulness of the
reduced projection map in cad construction (algorithms for cad
computation using the new projection are presented in Chapter 5).
Section 3 contains mathematical details that substantiate the results
in Section 2.

\section{The Original Cad Algorithm and Its Projection Map}

The purpose of this section is to provide a quick introduction to the
cad algorithm. More detailed accounts appear in a number of sources,
e.g.~[COL75], [ARN81], [ACM84a].

Let $A$ be a finite set of $r$-variate polynomials, $r \ge 1$. An
\emph{$A$-invariant cylindrical algebraic decomposition (cad) of $\R^r$}
partitions $\R^r$ into a finite collection of semialgebraic cells in
each of which every polynomial in $A$ is sign-invariant. A more precise
definition of cad is given in [ACM84a]. An $A$-invariant cad of the
plane, where
\[
A = \{ F(x, y) = y^4 - 2y^3 + y^2 - 3x^2 y + 2x^4 \},
\]
is depicted in Figure~3.1.1. Note that $F(x, y)$ vanishes in each of the
0-cells and ``curved'' 1-cells, and is either positive or negative
throughout each of the remaining cells.

\begin{center}
\textit{[Figure 3.1.1: illustration of an $A$-invariant cad of the plane
for the polynomial $F(x,y) = y^4 - 2y^3 + y^2 - 3x^2 y + 2x^4$, showing
0-cells, curved 1-cells along the zero locus of $F$, and the
complementary 2-cells.]}
\end{center}

The cad algorithm accepts as input a finite set $A$ of integral
polynomials in $r$ variables, and yields as output a description of an
$A$-invariant cad $D$ of $\R^r$. The description of $D$ takes the form
of a list of sample points and (if required) defining formulas for the
cells of $D$. The algorithm consists of three phases: \emph{projection}
(computing successive sets of polynomials in one fewer variables, the
zeros of each set containing a projection of the ``critical zeros'' in
the next higher dimensional space), \emph{base} (constructing a cad of
$\R^1$), and \emph{extension} (successive extension of the cad of
$\R^i$ to a cad of $\R^{i+1}$, $i = 1, 2, \ldots, r-1$). We shall
describe each of the phases in turn.

For a set $A$ of $r$-variate integral polynomials, we shall define a
set $\PROJ(A)$ of $(r-1)$-variate polynomials whose zeros contain the
projection of the ``critical locus'' of $A$. We shall show that for any
$\PROJ(A)$-invariant $j$-cell $R$ in $\R^{r-1}$, the cylinder over $R$
(i.e.\ $R \times \R$) is partitioned by the zeros of $A$ into a finite
number of pairwise disjoint $j$-cells (the graphs of continuous
functions defined on $R$) and their complementary $(j+1)$-cells. This
property will enable us to carry out the extension of a
$\PROJ(A)$-invariant cad of $\R^{r-1}$ to an $A$-invariant cad of
$\R^r$.

We introduce some terminology first. The \emph{reductum} of a polynomial
$F$ is the polynomial $F - \ldt(F)$, where $\ldt(F)$ is the leading term
of $F$, i.e.\ the term of highest degree. Let $\red^k(F)$ denote the
$k$-th reductum of $F$. Let $F(x)$ and $G(x)$ be nonzero polynomials
over an integral domain $D$, $F$ and $G$ not both constant. Let
$n = \min(\deg(F), \deg(G))$. For $0 \le j \le n$, let $S_j(F, G)$
denote the \emph{$j$-th subresultant polynomial} of $F$ and $G$, an
element of $D[x]$ of degree $\le j$. (Each coefficient of $S_j(F, G)$
is the determinant of a certain matrix of $F$ and $G$ coefficients;
see [LOO82b], [BRT71], or [COL75] for the exact definition. Subresultant
polynomials are intimately connected with polynomial remainder
sequences: the fundamental theorem of polynomial remainder sequences
[BRT71] makes the connections explicit. The fact concerning subresultant
polynomials of chief importance to us is that if $D$ is a unique
factorization domain, then where $k = \deg(\gcdop(F, G))$, we have
$S_j(F, G) = 0$ for $0 \le j < k$, and $S_k(F, G) = \gcdop(F, G)$ modulo
multiplication by elements of $D$.) For $0 \le j \le n$, the
\emph{$j$-th subresultant} of $F$ and $G$, written $\subres_j(F, G)$,
is the coefficient of $x^j$ in $S_j(F, G)$. Note that
$\subres_0(F, G) = \res(F, G)$. Suppose now that the characteristic of
$D$ is $0$. Let $\ldcf(F)$ denote the leading coefficient of $F$, and
assume that $(n :=) \deg(F) > 0$. Then $n \ldcf(F) \ne 0$, and so the
degree of $F'$, the derivative of $F$, is $n-1$. By considering the
Sylvester matrix of $F$ and $F'$, it can be seen that, for
$0 \le j \le n-1$, $\ldcf(F)$ is a divisor of $\subres_j(F, F')$. For
$0 \le j \le n-1$, we define the \emph{$j$-th subdiscriminant} of $F$,
written $\subdiscr_j(F)$, by the equation
\[
\ldcf(F) \cdot \subdiscr_j(F) = (-1)^{n(n-1)/2} \subres_j(F, F')
\]
(cf.\ definition of $\discr(F)$ in Sec.~2.3). Note that
$\subdiscr_0(F) = \discr(F)$.

Let $A$ be a set of integral polynomials in the $r$ variables
$x_1, \ldots, x_r$. We regard the elements of $A$ as polynomials in
$x_r$ over the ring of polynomials in $x_1, \ldots, x_{r-1}$. We now
define several sets of $(r-1)$-variate polynomials. The \emph{reducta
set of $A$}, written $\red(A)$, is the set
\[
\{ \red^k(F) : F \in A,\ 0 \le k \le \deg(F),\ \red^k(F) \ne 0 \}.
\]
The \emph{coefficient set of $A$}, written $\coeff(A)$, is the set
\[
\{ f(x_1, \ldots, x_{r-1}) : f \text{ is a non-zero coefficient of some }
F \in A \}.
\]
The \emph{subdiscriminant set of $A$}, written $\subdiscr(A)$, is
defined to be
\[
\{ \subdiscr_j(F) : F \in A,\ 0 \le j < \deg(F') \}.
\]
The \emph{subresultant set of $A$}, written $\subres(A)$, is defined to
be
\[
\{ \subres_j(F, G) : F, G \in A,\ F \ne G,\
0 \le j < \min(\deg(F), \deg(G)) \}.
\]
We can now define the projection of $A$. Suppose first that $A$ is a
\emph{squarefree basis} (that is, the elements of $A$ have positive
degree, and are primitive, squarefree and pairwise relatively prime).
Then we define $\PROJ(A)$ to be the set
\[
\coeff(A) \cup \subdiscr(\red(A)) \cup \subres(\red(A)).
\]
Now let $A$ be any set of $r$-variate integral polynomials. Let
$\cont(A)$ be the set of non-constant contents of elements of $A$, and
let $\prm(A)$ be the set of primitive parts of positive degree of
elements of $A$. Then where $B$ is the finest squarefree basis for
$\prm(A)$ (that is, the set of irreducible divisors of elements of
$\prm(A)$), we set
\[
\PROJ(A) = \cont(A) \cup \PROJ(B).
\]

The following concept is crucial:

\begin{definition}
A real polynomial $F(x, x_r)$ is said to be \emph{delineable} on a
subset $S$ of $\R^{r-1}$ if
\begin{enumerate}
\item[(1)] the portion of the real variety of $F$ lying in the cylinder
over $S$ consists of the union of the graphs of some $k \ge 0$
continuous functions $\theta_1 < \cdots < \theta_k$ from $S$ to $\R$;
and
\item[(2)] there exist integers $m_1, \ldots, m_k \ge 1$ such that for
every $\alpha \in S$, the multiplicity of the root $\theta_i(\alpha)$
of $F(\alpha, x_r)$ is $m_i$.
\end{enumerate}
\end{definition}

In the above definition, the graphs of the functions $\theta_i$ are
called the \emph{$F$-sections} over $S$. The regions between successive
$F$-sections are called \emph{$F$-sectors}.

We can now state the main theorem about $\PROJ$:

\begin{theorem*}[3.1.1]
Let $A$ be a finite set of integral polynomials in $r$ variables and
let $c \subseteq \R^{r-1}$ be a connected set in which every element of
$\PROJ(A)$ is sign-invariant. Then every element of $A$ is either
delineable or identically zero on $c$, and the sections of $A$ over $c$
are pairwise disjoint.
\end{theorem*}

We briefly discuss the ideas involved in the proof. It is shown in
[COL75] that for a polynomial $F(x, x_r)$ whose leading coefficient
does not vanish on a set $c \subseteq \R^{r-1}$, the invariance of the
number of distinct roots of $F(\alpha, x_r)$ as $\alpha$ varies within
$c$ implies the delineability of $F$ on $c$. On the other hand, the
number of distinct roots of $F(\alpha, x_r)$ is equal to $d - k$, where
$d$ is the degree of $F$ in $x_r$ and $k$ is the least integer
$j \ge 0$ such that $\subdiscr_j(F)(\alpha) \ne 0$ (this follows from
the fundamental theorem of polynomial remainder sequences [BRT71]).

Let us assume the hypotheses of the main theorem, and assume first that
$A$ is a squarefree basis. Let $F$ be an element of $A$ which is not
identically zero on $c$. Then as the elements of $\coeff(A)$ are
sign-invariant on $c$, $F$ is identical with some reductum $G$ of $F$
on $c$, such that $G$ has non-vanishing leading coefficient on $c$. As
the elements of $\subdiscr(\red(A))$ are sign-invariant on $c$,
$\subdiscr_j(G)$ is sign-invariant on $c$, for $0 \le j \le \deg(G')$.
Thus, by the above remarks, the number of distinct roots of
$G(\alpha, x_r)$ is invariant throughout $c$, hence $G$ (and therefore
$F$) is delineable on $c$. The invariance of $\subres(\red(A))$ on $c$
ensures that the various sections of elements of $A$ on $c$ are
pairwise disjoint.

If $A$ is an arbitrary set of polynomials, then we can apply the above
argument to the basis $B$ for $\prm(A)$. The required conclusions for
$A$ follow immediately. The reader can consult either [COL75] or
[ACM84a] for a more detailed proof. $\square$

Let $A$ be the input to the cad algorithm. Let us assume for the
present that defining formula construction is not required. In the
projection phase of the algorithm we compute $\PROJ(A)$,
$\PROJ(\PROJ(A)) = \PROJ^2(A)$, and so on, until we compute
$\PROJ^{r-1}(A)$. It is the task of the base phase to construct a
$\PROJ^{r-1}(A)$-invariant cad of $\R^1$. We begin the base phase by
constructing the set of all distinct irreducible factors of nonzero
elements of $\PROJ^{r-1}(A)$. We then isolate the real roots of (the
product of) these factors. We thus have an exact representation for
each root consisting of its minimal polynomial and an isolating interval
(see [LOO82a], Section 1 for more information on this). The cad of
$\R^1$ consists of the collection of these roots (the 0-cells)
together with the collection of complementary open intervals (the
1-cells). Sample points for the 0-cells are the roots themselves
represented as above, while rational sample points for the 1-cells can
be readily chosen. Defining formulas for the 0-cells and 1-cells can be
obtained in a straightforward manner (see p.~45 of [ARN81]).

Let $D^*$ be the cad of $\R^1$ constructed as above. Let us first
consider the extension of $D^*$ to a cad of $\R^2$. In the projection
phase we computed a set $J = \PROJ^{r-2}(A)$. Let $c$ be a cell of
$D^*$. Then by Theorem~3.1.1, the cylinder over $c$ is partitioned by
the zeros of $J$ into a finite number of sections and sectors. These
sections and sectors will be cells of our decomposition of $\R^2$. We
now describe how to determine the number of these cells over $c$, as
well as a sample point for each cell. Let $\alpha$ be the sample point
for $c$, and let $J_\alpha$ be the product of all elements $G$ of $J$
for which $G(\alpha, x_2) \ne 0$. Using algorithms for exact arithmetic
in $\Q(\alpha)$ [LOO82a], we construct $J_\alpha(\alpha, x_2)$. We
isolate the real roots of $J_\alpha(\alpha, x_2)$, which provides us
with the number of sections and sectors of $J$ lying over $c$. Now
$(\alpha, \beta)$ lies on a section of $J$ over $c$ if and only if
$\beta$ is a root of $J_\alpha(\alpha, x_2)$. For each such $\beta$, we
use the representation for $\alpha$, the isolating interval for
$\beta$, and the algorithms \textsc{NORMAL} and \textsc{SIMPLE} of
[LOO82a] to construct a primitive element $\gamma$ for
$\Q(\alpha, \beta)$; we use $\gamma$ to construct an appropriate
representation for the algebraic point $(\alpha, \beta)$, by expressing
$\alpha$ and $\beta$ as polynomials in $\gamma$. We readily get sector
sample points of the form $(\alpha, r)$, with $r$ rational. After
processing each cell $c$ of $D^*$ in this way, we have determined a
cad of $\R^2$ and constructed a sample point for each cell.

Extension from $\R^{i-1}$ to $\R^i$ for $3 \le i \le r$ is essentially
the same as from $\R^1$ to $\R^2$. The only difference is that a sample
point in $\R^{i-1}$ has $i-1$, instead of just one, coordinates. But
where $\alpha$ is the primitive element of an $\R^{i-1}$ sample point,
and $F$ is a polynomial in $r$ variables, arithmetic in $\Q(\alpha)$
still suffices for constructing the univariate polynomial over
$\Q(\alpha)$ that results from substituting the coordinates
$(\alpha_1, \ldots, \alpha_{i-1})$ for $(x_1, \ldots, x_{i-1})$ in $F$.

We have also to define the \emph{augmented projection} $\APROJ(A)$ of
$A$, which is used when defining formulas are to be constructed. First
some notation:

\begin{definition}
Let us take the degree of the zero polynomial to be $-1$. Let
\[
f(x, x_r) = f_0(x) x_r^d + f_1(x) x_r^{d-1} + \cdots + f_d(x)
\]
be a polynomial in $\R[x, x_r]$ and let $S$ be a subset of $\R^{r-1}$.
Say that $f$ is \emph{degree-invariant in $S$} if there exists $l$,
$-1 \le l \le d$, called the \emph{degree of $f$ in $S$} (written
$\deg_S(f)$), such that for every point $p$ of $S$, the univariate
polynomial $f(p, x_r)$ has degree $l$. If $f$ is degree-invariant and
of non-negative degree $l$ in $S$, then
$f_0 = f_1 = \cdots = f_{d-l-1} = 0$ in $S$, while $f_{d-l} \ne 0$ in
$S$ (the converse is also true).
\end{definition}

Let $F^{(j)}$ denote the $j$-th derivative of $F$. Define the
\emph{derivative set} $\der(A)$ of $A$ to be the set
\[
\{ F^{(j)} : F \in A,\ 0 \le j < \deg(F) \}.
\]
Suppose first that $A$ is a squarefree basis. Then we define
$\APROJ(A)$ to be the set
\[
\subdiscr(\der(\red(A))) \cup \PROJ(A).
\]
If $A$ is an arbitrary set of polynomials then where $B$ is the finest
squarefree basis for $\prm(A)$, we set
\[
\APROJ(A) = \cont(A) \cup \APROJ(B).
\]

The role of the augmented projection in defining formula construction
is indicated by the following two theorems:

\begin{theorem*}[3.1.2]
Let $A$ be a finite set of integral polynomials in $r$ variables and
let $c \subseteq \R^{r-1}$ be a connected set in which every element of
$\APROJ(A)$ is sign-invariant. Then for every $F \in A$, $F$ is
degree-invariant on $c$, and $F^{(j)}$ is delineable on $c$, for every
$j$, $0 \le j \le \deg_c(F)$.
\end{theorem*}

\begin{theorem*}[3.1.3]
Let $F$ be an integral polynomial in $r$ variables. Let
$c \subseteq \R^{r-1}$ be a connected set on which $F$ is
degree-invariant and not identically zero, and on which $F^{(j)}$ is
delineable for $0 \le j \le \deg_c(F)$. Then, given a sample point and
a defining formula for $c$, we can construct defining formulas for the
$F$-sections and $F$-sectors over $c$.
\end{theorem*}

For the proofs of these theorems one should consult Sec.~2.5 of
[ARN81]. The following abstract algorithm indicates how augmented
projections are used in the cad algorithm, and generally serves to
summarize the discussion of this section.

\bigskip
\noindent \textbf{CAD}$(r, A, k; S, F)$

\medskip
[Cylindrical algebraic decomposition. $A$ is a set of integral
polynomials in $r$ variables, $r \ge 1$. $k$ satisfies $0 \le k \le r$.
$S$ is a list of sample points for an $A$-invariant cad $D$ of $\R^r$.
If $k \ge 1$, $F$ is a list of defining formulas for the induced cad of
$\R^k$, and if $k = 0$, $F$ is the empty list.]

\begin{enumerate}
\item[(1)] [Initialize.] Set $B \leftarrow$ the finest squarefree basis
for $\prm(A)$. Set $S \leftarrow ()$ and $F \leftarrow ()$.

\item[(2)] [$r = 1$.] If $r > 1$ then go to (3). Isolate the real roots
of $B$. Construct sample points for the cells of $D$ and add them to
$S$. If $k = 1$, then construct defining formulas for the cells of $D$
and add them to $F$. Exit.

\item[(3)] [$r > 1$.] If $k < r$ then set $P \leftarrow \PROJ(A)$ and
$k' \leftarrow k$; otherwise set $P \leftarrow \APROJ(A)$ and
$k' \leftarrow k - 1$. Call CAD recursively with inputs $r-1$, $P$, and
$k'$ to obtain outputs $S'$ and $F'$ which specify a $P$-invariant cad
$D'$ of $\R^{r-1}$. For each cell $c$ of $D'$, let $\alpha$ denote the
sample point for $c$, and carry out the following sequence of steps:
set $B^* \leftarrow$ the set of all $B_i(\alpha, x_r)$ such that
$B_i \in B$ and $B_i(\alpha, x_r) \ne 0$; isolate the real roots of
$B^*$; use $\alpha$ and the isolating intervals for the roots of $B^*$
to construct sample points for the $B$-sections and $B$-sectors over
$c$, adding them to $S$; if $k = r$, then, from the defining formula
for $c$, construct defining formulas for the $B$-sections and
$B$-sectors over $c$, adding them to $F$, and if $k < r$, set
$F \leftarrow F'$. Exit. $\square$
\end{enumerate}

\section{A Reduced Projection Map and New Projection Theorem}

In this section we introduce a new projection map $P$ which is a
`reduced' version of the map $\PROJ$ discussed in Section~3.1. We state
a theorem analogous to Theorem~3.1.1 which implies that under certain
assumptions on a given set $A$ of polynomials, any cad of $\R^{r-1}$
such that every element of $P(A)$ has constant \emph{order} (as
opposed to sign~--- recall the definition of order given in
Section~2.3) throughout every cell, can be lifted (or extended) to a
cad of $\R^r$ such that every element of $A$ has constant order
throughout every cell. The mathematical result underlying this
analogue of Theorem~3.1.1 will be termed the \emph{lifting theorem},
and is the main contribution of this thesis. In this section we state
the lifting theorem and derive the important consequences for cad
construction from it. We postpone the proof of the lifting theorem to
the next section of this chapter. In Chapter~5, we present algorithms
for cad construction based upon the lifting theorem.

Before we can state the lifting theorem we need to make a few
definitions.

\begin{definition}
Let $K = \R$ or $\C$. Let $U$ be an open subset of $K^r$ and let
$f: U \to K$ be an analytic function. We say that $f$ is
\emph{order-invariant in} a subset $S$ of $U$ provided that the order
of $f$ is the same at every point of $S$.
\end{definition}

\begin{example}
Let $f: \R^2 \to \R$ be given by $f(x, y) = x^2 - y^2$. Let $C_f$ be
the curve defined by $f(x, y) = 0$ and let $S = C_f - \{0\}$. Then $f$
is order-invariant in $S$ ($\ord_{(x_0, y_0)} f = 1$ for every point
$(x_0, y_0)$ of $S$). However $f$ is not order-invariant in $C_f$
($\ord_0 f = 2$).
\end{example}

\begin{remark}
Let $U$ and $V$ be open subsets of $K^n$, let $G: U \to V$ be an
analytic isomorphism, and let $f: V \to K$ be an analytic function. By
Theorem~2.1.5 the composite function $f \circ G$ is analytic in $U$.
It follows from Theorem~2.3.1 that, for every point $p$ of $U$,
\[
\ord_{G(p)} f = \ord_p f \circ G.
\]
Hence, if $S$ is a subset of $U$, then $f$ is order-invariant in
$G(S)$ if and only if $f \circ G$ is order-invariant in $S$.
\end{remark}

Throughout the remainder of this section and the next, let $x$ denote
the $(r-1)$-tuple $(x_1, \ldots, x_{r-1})$ and let $(x, x_r)$ denote
$(x_1, \ldots, x_{r-1}, x_r)$.

We now introduce the concept of analytic delineability.

\begin{definition}
Let $S$ be a connected submanifold of $\R^{r-1}$. We saw in Sec.~2.2
that one can define the notion of an analytic function from $S$ into
$\R$. If $f(x, x_r)$ is a real polynomial then we say $f$ is
\emph{analytic delineable} on $S$ if $f$ is delineable on $S$ by means
of some $k \ge 0$ analytic functions $\theta_1 < \cdots < \theta_k$
from $S$ into $\R$.
\end{definition}

\begin{remark}
Suppose that $f(x, x_r)$ is analytic delineable on the connected
$s$-submanifold $S$ of $\R^{r-1}$. Then each $f$-section over $S$ is a
connected $s$-submanifold of $\R^r$ (Theorem~2.2.3) and each $f$-sector
over $S$ is a connected $(s+1)$-submanifold of $\R^r$ (Theorem~2.2.4).
\end{remark}

We can now state our main mathematical result.

\begin{theorem*}[3.2.1 (Lifting theorem)]
Let $f(x, x_r)$ be a squarefree polynomial in $\R[x, x_r]$ of positive
degree in $x_r$ and with non-zero discriminant $D(x)$. Let $S$ be a
connected submanifold of $\R^{r-1}$ in which $f$ is degree-invariant
and of non-negative degree, and in which $D$ is order-invariant. Then
$f$ is analytic delineable on $S$ and is order-invariant in each
$f$-section over $S$.
\end{theorem*}

The proof of this theorem is quite long, and is supplied in the next
section of this chapter. The theorem will enable us to show that our
reduced projection map $P$, to be defined shortly, is essentially as
good (for the purpose of cad construction) as its larger counterpart
$\PROJ$.

We now proceed to define the map $P$. Let $A$ be a set of $r$-variate
integral polynomials. Recall the definition of the coefficient set
$\coeff(A)$ from the previous section. We define the \emph{discriminant
set of $A$}, written $\discr(A)$, to be the set
\[
\{ \discr(F) : F \in A,\ \deg(F) \ge 2 \},
\]
and the \emph{resultant set of $A$}, written $\res(A)$, to be the set
\[
\{ \res(F, G) : F, G \in A,\ F \ne G,\ \deg(F), \deg(G) \ge 1 \}.
\]
We can now define $P$. Suppose first that $A$ is a squarefree basis.
We set
\[
P(A) = \coeff(A) \cup \discr(A) \cup \res(A).
\]
Now let $A$ be any set of $r$-variate polynomials. We set
\[
P(A) = \cont(A) \cup P(B),
\]
where $B$ is the finest squarefree basis for $\prm(A)$.

Noting that the resultant of two polynomials $F$ and $G$ is equal to
$\subres_0(F, G)$, and that the discriminant of $F$ is equal to
$\subdiscr_0(F)$, we can compare the definitions of $\PROJ$ and $P$.
It is clear that $P(A)$ is a subset of $\PROJ(A)$. There are two
respects in which $\PROJ$ and $P$ differ. The first is that $\PROJ$
forms resultants, etc.\ of elements of the reductum set of $B$, whereas
$P$ forms resultants, etc.\ from just $B$ itself. The second is that
$\PROJ$ constructs full sequences of $\subdiscr_j$'s and
$\subres_j$'s, whereas $P$ constructs only the 0-th order $\subdiscr$'s
and $\subres$'s (i.e.\ discriminants and resultants only).

Before deriving a counterpart of Theorem~3.1.1, we give a useful lemma.

\begin{lemma*}[3.2.2]
Let $K = \R$ or $\C$. Let $F_1, \ldots, F_n$ be nonconstant analytic
functions defined in an open subset $U$ of $K^r$, and let
$f = F_1 \cdots F_n$. Let $S$ be a connected subset of $U$. Then $f$
is order-invariant in $S$ if and only if each $F_i$ is order-invariant
in $S$.
\end{lemma*}

The proof is straightforward. $\square$

We now state our central result pertaining to $P$:

\begin{theorem*}[3.2.3]
Let $A$ be a finite squarefree basis consisting of $r$-variate integral
polynomials, $r \ge 2$, and let $S$ be a connected submanifold of
$\R^{r-1}$. Suppose that each element of $A$ is not identically zero
on $S$, and that each element of $P(A)$ is order-invariant in $S$.
Then each element of $A$ is degree-invariant and analytic-delineable
on $S$, the sections of $A$ over $S$ are pairwise disjoint, and each
element of $A$ is order-invariant in every section of $A$ over $S$.
\end{theorem*}

\begin{proof}
Assume that some element of $A$ is nonconstant (otherwise the theorem
is trivial) and let $F_1, \ldots, F_n$ be the nonconstant elements of
$A$. Let $f = \prod_{i=1}^{n} F_i$. Then $f$ is squarefree, as the
$F_i$ are squarefree and pairwise relatively prime. Let $D(x)$ be the
discriminant of $f(x, x_r)$. By Theorem~2.3.3,
\[
D = \prod_{i=1}^{n} \discr(F_i) \cdot \prod_{i < j} \res(F_i, F_j)^2.
\]
By hypothesis, each $\discr(F_i)$ and each $\res(F_i, F_j)$ is
order-invariant in $S$. Hence, by Lemma~3.2.2, $D$ is order-invariant
in $S$. Furthermore, $f$ is not identically zero on $S$, and is
degree-invariant on $S$. Hence, by the lifting theorem, $f$ is analytic
delineable on $S$, and is order-invariant in each $f$-section over $S$.
By Lemma~3.2.2, each $F_i$ is order-invariant in every such section.
The conclusions of the theorem now follow.
\end{proof}

We shall also define a reduced augmented projection map $AP$, and
establish a counterpart of Theorem~3.1.2. Let $A$ be a squarefree
basis of integral polynomials in $\Z[x, x_r]$. We define a slight
variant of the derivative set of $A$: where $\gsfd(P)$ denotes the
greatest squarefree divisor of the polynomial $P$, let
\[
\der^*(A) = \{ \gsfd(\prm(F^{(j)})) : F \in A,\ 0 < j < \deg(F) \}.
\]
We set
\[
AP(A) = \discr(\der^*(A)) \cup P(A).
\]
For an arbitrary set $A$ we set
\[
AP(A) = \cont(A) \cup AP(B),
\]
where $B$ is the finest squarefree basis for $\prm(A)$.

The following result is our counterpart of Theorem~3.1.2:

\begin{theorem*}[3.2.4]
Let $A$ be a finite squarefree basis of integral polynomials in $r$
variables, where $r \ge 2$, and let $S$ be a connected submanifold of
$\R^{r-1}$. Suppose that each element of $A$ is not identically zero
on $S$, and that each element of $AP(A)$ is order-invariant in $S$.
Then, for each $F \in A$, $F$ is degree-invariant on $S$, and
$F^{(j)}$ is delineable on $S$, for every $j$, $0 \le j \le \deg_S(F)$.
\end{theorem*}

\begin{proof}
Let $F \in A$. The degree-invariance of $F$ on $S$ is obvious. Let
$0 \le j \le \deg_S(F)$. Then $F^{(j)}$ is degree-invariant and not
identically zero on $S$ (as $F$ is). Let $G$ be the greatest squarefree
divisor of the primitive part of $F^{(j)}$. Then, as $G$ is a divisor
of $F^{(j)}$, $G$ is degree-invariant and not identically zero on $S$.
By hypothesis, the discriminant of $G$ is order-invariant in $S$.
Hence, by the lifting theorem, $G$ is delineable on $S$. As the content
of $F^{(j)}$ is non-zero in $S$, it follows that $F^{(j)}$ is
delineable on $S$.
\end{proof}

\section{Proof of the Lifting Theorem}

The lifting theorem can be regarded as an adaptation to Euclidean space
of certain aspects of O.\ Zariski's work on equisingularity over the
complex field. In this section we give a straightforward proof of the
lifting theorem based upon a recent result due to Zariski [ZAR75]. We
do not, however, require the reader to consult Zariski's original paper
and to attempt to see how the result stated therein is applicable.
Instead, we provide in Chapter~4 a self-contained, and in some respects
original, rendition of the relevant theorems due to Zariski.

The proof presented in this section is quite detailed, and hence
rather long. The main thread of the proof can be gleaned from the
first few pages.

\begin{notation}
We denote by $Z(S)$ the cylinder over a subset $S$ of $\R^{r-1}$ (i.e.\
$Z(S) = S \times \R$).
\end{notation}

\begin{proof}[Proof of 3.2.1]
Let the dimension of $S$ be $s$, where $0 \le s \le r-1$. If $s = 0$
then the theorem is trivial. For $1 \le s \le r-1$, the theorem follows
from the following assertion as to the ``local delineability'' of $f$:

\begin{assertion}[1]
For each point $p$ of $S$ there is a neighborhood $N \subseteq \R^{r-1}$
of $p$ such that $f$ is analytic delineable on $S \cap N$ and $f$ is
order-invariant in each $f$-section over $S \cap N$.
\end{assertion}

We explain now how the theorem can be deduced from the above assertion.
Let Assertion~1 hold. Then it follows by connectedness of $S$ that the
number, say $k \ge 0$, of distinct real roots of the univariate
polynomial $f(p, x_r)$ does not depend upon the particular point $p$ of
$S$ (that this number is locally constant is an immediate consequence
of Assertion~1). If $k = 0$ then the theorem is vacuously true, so let
$k \ge 1$. Then one can define $k$ functions $\Theta_1, \ldots,
\Theta_k$ from $S$ into $\R$ by setting $\Theta_i(p)$ equal to the
$i$-th real root of $f(p, x_r)$, for $p$ in $S$ and $1 \le i \le k$.
Clearly the graphs of the $\Theta_i$ comprise the portion of the real
variety of $f$ lying in the cylinder over $S$. Let $p$ be a point of
$S$. Then by Assertion~1 there is a neighborhood $N \subseteq \R^{r-1}$
of $p$ and $k$ analytic functions from $S \cap N$ into $\R$, say
$\theta_1 < \cdots < \theta_k$, whose graphs comprise the portion of
the real variety of $f$ which lies in $Z(S \cap N)$, and such that $f$
is order-invariant in the graph of each $\theta_i$. Let
$1 \le i \le k$. As $\Theta_i(x) = \theta_i(x)$ for all $x$ in
$S \cap N$, it follows that $\Theta_i$ is analytic near $p$, and hence
in $S$. Thus $f$ is analytic delineable on $S$. By connectedness of
$S$, $f$ is order-invariant in the graph of each $\Theta_i$. The
conclusions of the lifting theorem have been verified.

It remains to establish Assertion~1 (i.e., the local delineability of
$f$). Let $p$ be a point of $S$, and let the degree of $f(p, x_r)$ be
$l$. Then $l \ge 0$ by hypothesis. If $l = 0$ then Assertion~1 follows
immediately by the degree-invariance of $f$ in $S$. Therefore, let us
henceforth assume that $l > 0$. Let $\alpha_1 < \cdots < \alpha_k$,
$k \ge 0$, be the distinct real roots of $f(p, x_r)$, let
$\alpha_{k+1}, \ldots, \alpha_t$, $k \le t$, be the distinct non-real
roots of $f(p, x_r)$, and let $m_i$ be the multiplicity of the root
$\alpha_i$, for $1 \le i \le t$. Then $\sum_{i=1}^{t} m_i = l$. Let
$\kappa$ be the minimum separation between the roots of $f(p, x_r)$:
that is, $\kappa = \min\{ |\alpha_i - \alpha_j| : 1 \le i < j \le t \}$.
Let $0 < \epsilon < \kappa/3$ and let $C_i$ be the circle in the
complex plane of radius $\epsilon$ centred at $\alpha_i$. By the root
continuity principle (Theorem~2.3.2) there is a neighborhood
$N_0 \subseteq \R^{r-1}$ of $p$ such that for every $x \in N_0$, each
$C_i$ contains exactly $m_i$ roots (multiplicities counted) of
$f(x, x_r)$. Consider the following assertion, which implies that if
$x$ is sufficiently close to $p$, and belongs to $S$, then all of the
$m_i$ roots of $f(x, x_r)$ inside $C_i$ are coincident:

\begin{assertion}[2]
For each $i$, $1 \le i \le k$, there is a neighborhood
$N_i \subseteq N_0$ of $p$ and an analytic function $\theta_i$ from
$S \cap N_i$ into $(\alpha_i - \epsilon, \alpha_i + \epsilon)$ such
that for every $x \in S \cap N_i$ and every
$\alpha \in (\alpha_i - \epsilon, \alpha_i + \epsilon)$,
\begin{equation}
f(x, \alpha) = 0 \text{ if and only if } \alpha = \theta_i(x), \tag{3.3.1}
\end{equation}
and such that $f$ is order-invariant in the graph of $\theta_i$.
\end{assertion}

We say that the above statement asserts the \emph{nonsplitting} of the
roots $\alpha_i$, as $x$ varies within $S$ near $p$.

We shall now verify Assertion~1 on the strength of Assertion~2. Assume
that Assertion~2 holds. Let $N = \bigcap_{i=0}^{k} N_i$, and consider
the domain of each $\theta_i$ to be $S \cap N$. We shall show that the
portion of the real variety of $f$ (denoted by $V(f)$) which lies in
the cylinder over $S \cap N$ is the union of the graphs of the
$\theta_i$: that is,
\begin{equation}
V(f) \cap Z(S \cap N) = \bigcup_{i=1}^{k} \operatorname{graph}(\theta_i), \tag{3.3.2}
\end{equation}
(where the right-hand side of the equation is understood to be the
empty set in case $k = 0$). Let $(x, \alpha)$ be an element of the
right-hand side of (3.3.2). Then $x \in S \cap N$ and
$\alpha = \theta_i(x)$ for some $i$, $1 \le i \le k$. Hence, as the
range of $\theta_i$ is $(\alpha_i - \epsilon, \alpha_i + \epsilon)$,
$\alpha \in (\alpha_i - \epsilon, \alpha_i + \epsilon)$. Therefore, by
(3.3.1), $f(x, \alpha) = 0$, and so $(x, \alpha)$ is contained in the
left-hand side of (3.3.2). Conversely, let $(x, \alpha)$ be an element
of $V(f) \cap Z(S \cap N)$. Then $x \in S \cap N$, $\alpha \in \R$, and
$f(x, \alpha) = 0$. As $f$ is degree-invariant in $S$ and the degree
of $f(p, x_r)$ is $l$, the degree of $f(x, x_r)$ is also $l$. As
$x \in N_0$, the interior of each $C_i$, $1 \le i \le t$, contains
exactly $m_i$ roots (multiplicities counted) of $f(x, x_r)$. Since
$\sum_{i=1}^{t} m_i = l$, every root of $f(x, x_r)$ is contained within
one of the $C_i$. Each $C_i$ with $k+1 \le i \le t$ contains no real
points, however, as the non-real roots of $f(p, x_r)$ occur in
conjugate pairs. Hence, as $\alpha$ is a real root of $f(x, x_r)$,
$\alpha$ must lie inside a $C_i$ with $1 \le i \le k$. Hence, by
(3.3.1), $\alpha = \theta_i(x)$, so $(x, \alpha)$ is contained in the
graph of $\theta_i$. We have shown that (3.3.2) holds. This completes
our proof of Assertion~1, on the strength of Assertion~2.

It therefore remains to establish Assertion~2 (i.e., the nonsplitting
of the roots $\alpha_i$ as $x$ varies within $S$ near $p$). We now
give a point-by-point summary of the steps involved in proving the
nonsplitting of the roots.

\begin{enumerate}
\item We fix on a particular $\alpha_i$, and observe that there is no
loss of generality in assuming that $\alpha_i = 0$. We are able to
dispose quite quickly of the case in which the dimension $s$ of $S$ is
equal to $r-1$. In this case, $S$ is just an open subset of
$\R^{r-1}$, and so both the discriminant of $f$ and the leading
coefficient of $f$ are nonzero in $S$. Hence, the nonsplitting of the
roots is obvious (all roots are simple), and the analytic function
$\theta_i$ of Assertion~2 is given by the implicit function theorem.

\item There remains the harder case $1 \le s \le r-2$. We choose
coordinates $(y_1, \ldots, y_{r-1})$ about the point $p$, such that
$S$ is defined locally by the equations $y_{s+1} = 0, \ldots,
y_{r-1} = 0$ in the new coordinate system (by Theorem~2.2.1). Let
$g(y, y_r)$ denote the function $f(x, x_r)$ transformed into the new
coordinates. Then $g(y, y_r)$ is a polynomial in $y_r$ whose
coefficients are analytic functions of $y$, defined near $0$ (the
analyticity here comes from the analyticity of the coordinate change
map). The discriminant $E(y)$ of $g(y, y_r)$ is analytic near $0$, and
is order-invariant in the linear subspace $T$ defined by
$y_{s+1} = 0, \ldots, y_{r-1} = 0$, near $0$ (as $D(x)$ is
order-invariant in $S$).

\item Each coefficient of $g(y, y_r)$ can be expanded in a convergent
power series about $0$ (by definition of analyticity). Each such power
series can be regarded as a complex power series in the complex
variables $(z_1, \ldots, z_{r-1}) = z$, where $z_i = y_i + i v_i$,
say. By the multivariate analogue (Theorem~2.1.2) of a well-known
result on convergence, each of these power series is absolutely
convergent in a neighborhood of $0$ in complex $(r-1)$-space
$\C^{r-1}$. In this way, each coefficient of $g(y, y_r)$, and hence
also $g(y, y_r)$, can be extended (uniquely) to complex space. We
write $g(z, z_r)$ to denote this extension of $g(y, y_r)$:
$g(z, z_r)$ is a pseudopolynomial near $0$ (see Section~2.1 for
definition). We observe that the discriminant $E(z)$ of $g(z, z_r)$
is order-invariant in the complex linear subspace $T^*$ defined by
$z_{s+1} = 0, \ldots, z_{r-1} = 0$, near $0$.

\item The next step is to focus our attention on the structure of the
complex variety $g(z, z_r) = 0$ in a neighborhood of the origin. We use
the Weierstrass Preparation Theorem (Theorem~2.1.10) from the theory
of several complex variables to do this. This theorem gives us a monic
pseudopolynomial $h(z, z_r)$ which has the same locus as $g(z, z_r)$
near the origin. It is shown that the discriminant $F(z)$ of
$h(z, z_r)$ is order-invariant in $T^*$, near $0$.

\item We can now apply Zariski's theorem (Theorem~4.1.1) to $h$. This
theorem implies the nonsplitting of the roots of $h(z, z_r)$, as $z$
varies within $T^*$, near $0$. In fact, the theorem yields an analytic
function $\psi(z_1, \ldots, z_s)$ defined near $0$ which gives the
unique distinct root $\alpha$ of $h(z, z_r)$, for
$z = (z_1, \ldots, z_s, 0, \ldots, 0)$ in $T^*$, sufficiently close to
$0$. The theorem also asserts order-invariance of $h$ along the locus
of this root.

\item All that remains is to retrace our path. The nonsplitting of the
roots of $g(y, y_r)$, as $y$ varies within $T$ near $0$, is immediate
(as $T \subseteq T^*$). Taking the restriction of $\psi$ to a
neighborhood of $0$ in real $s$-space $\R^s$ gives a real-valued
analytic function $\eta(y_1, \ldots, y_s)$ which represents the $i$-th
real root of $g(y, y_r)$ for $y$ in $T$, near $0$. Finally,
transforming back to the $(x, x_r)$-variables, we deduce the
nonsplitting of the roots of $f(x, x_r)$ as $x$ varies within $S$,
near $p$. We obtain an analytic function $\theta_i$ defined in $S$
near $p$, which represents the $i$-th real root of $f(x, x_r)$, for
$x$ in $S$, near $p$. Order-invariance is preserved under the
restriction and detransformation, and thus $f$ is order-invariant in
the graph of $\theta_i$.
\end{enumerate}

This concludes our summary of the proof of Assertion~2 (the
nonsplitting of the roots).

We now fill in the details of the proof of Assertion~2 outlined
above. Assume $k > 0$ ($k$ the number of real roots of $f(p, x_r)$)
and let $1 \le i \le k$. There is no loss of generality in assuming
that $\alpha_i = 0$. (For if $\alpha_i \ne 0$ then we can use a
translation of the $x_r$-coordinate to obtain a polynomial
$f'(x, x_r') = f(x, x_r' + \alpha_i)$ such that (a) $f'(p, x_r')$ has
a root of multiplicity $m_i$ at $x_r' = 0$ and no other roots either
inside or on the circle $C$ of radius $\epsilon$ about $0$; and (b)
the discriminant of $f'(x, x_r')$ is the same as that of $f(x, x_r)$.
Assuming that the required assertion has been proved for
$\alpha_i = 0$, we obtain a neighborhood $N_i$ of $p$ and an analytic
function $\theta_i' : S \cap N_i \to (-\epsilon, +\epsilon)$
corresponding to $f'$. We need simply ``de-translate'' $\theta_i'$
to find the desired
\[
\theta_i : S \cap N_i \to (\alpha_i - \epsilon, \alpha_i + \epsilon)
\]
corresponding to $f$.) Let $m = m_i$ and $C = C_i$: we seek a
neighborhood $N = N_i$ of $p$ and an analytic function
$\theta = \theta_i$ satisfying the required properties.

Let us first dispose of the

\medskip
\noindent \textbf{Case $s = r-1$.}
\medskip

By definition of submanifold there exists a neighborhood
$U \subseteq N_0$ of $p$ such that $U \subseteq S$. Let $f_0(x)$ be the
leading coefficient of $f(x, x_r)$, and let $R(x)$ be the resultant of
$f(x, x_r)$ and $\dfrac{\partial f}{\partial x_r}$. Recall that $D(x)$,
the discriminant of $f(x, x_r)$, is defined by the equation
\[
(-1)^{m(m-1)/2} R(x) = f_0(x) D(x).
\]
As $f$ is degree-invariant and $D$ is order-invariant in $S$,
$f_0 \ne 0$ and $D \ne 0$ in $U$ (since $f_0(x)$ and $D(x)$ are
non-zero polynomials). Hence $R \ne 0$ in $U$ and in particular
$R(p) \ne 0$. As $f_0(p) \ne 0$, $R(p)$ is equal to the resultant of
the univariate polynomials $f(p, x_r)$ and
$\dfrac{\partial f}{\partial x_r}(p, x_r)$. (Although this may seem
like a restatement of the definition of resultant, it actually depends
crucially on the fact that $f_0(p) \ne 0$. Recall that $R(x)$ is
defined to be a certain determinant in the coefficients of $f(x, x_r)$
and $\dfrac{\partial f}{\partial x_r}$. The only non-zero entries in
the first column of this determinant are $f_0(x)$ and $d f_0(x)$,
where $d$ is the degree of $f(x, x_r)$ (with respect to $x_r$). Hence
if $f_0(q) = 0$ then $R(q) = 0$. However the resultant of the
univariate polynomials $f(q, x_r)$ and $\dfrac{\partial f}{\partial x_r}
(q, x_r)$ could be non-zero in spite of $f_0(q) = 0$.) Since
$f(p, 0) = 0$, we must therefore have
$\dfrac{\partial f}{\partial x_r}(p, 0) \ne 0$. Hence, by the implicit
function theorem (obtained by taking $s = r-1$ in Theorem~2.1.7),
there are neighborhoods $N' \subseteq U$ of $p$ and
$W \subseteq (-\epsilon, +\epsilon)$ of $0$, and an analytic function
$\theta : S \cap N' = N' \to \R$, such that for all $x \in N'$ and all
$\alpha \in W$,
\begin{equation}
f(x, \alpha) = 0 \text{ if and only if } \alpha = \theta(x). \tag{3.3.3}
\end{equation}
By continuity of $\theta$, $N'$ may be refined to a neighborhood
$N \subseteq N'$ of $p$ such that $\theta(N) \subseteq W$. Let
$x \in S \cap N = N$ and let $\alpha \in (-\epsilon, +\epsilon)$. We
shall check that (3.3.1) holds. Assume that $\alpha = \theta(x)$.
Then $\alpha \in W$, so by (3.3.3), $f(x, \alpha) = 0$. Assume
conversely that $f(x, \alpha) = 0$. Now $m = 1$ as
$\dfrac{\partial f}{\partial x_r}(p, 0) \ne 0$. Hence we must have
$\alpha = \theta(x)$ (otherwise $f(x, x_r)$ would have at least two
distinct roots, $\alpha$ and $\theta(x)$, inside $C$, contradicting
$m = 1$). So (3.3.1) has been verified. Let $x \in N$ and let
$\alpha = \theta(x)$. As $m = 1$, $x_r = \alpha$ is a simple root of
$f(x, x_r)$. Thus $\dfrac{\partial f}{\partial x_r}(x, \alpha) \ne 0$,
and so $\ord_{(x, \alpha)} f = 1$. Hence $f$ is order-invariant in the
graph of $\theta$. Assertion~2 has been established.

There remains the

\medskip
\noindent \textbf{Case $1 \le s \le r-2$.}
\medskip

It will be convenient for us to use a coordinate system about $p$ with
respect to which $S$ is locally an $s$-dimensional linear subspace of
$\R^{r-1}$. By Theorem~2.2.1 there exists a neighborhood
$U \subseteq N_0$ of $p$ and a coordinate system $\Phi : U \to V$
about $p$, $\Phi = (\phi_1, \ldots, \phi_{r-1})$, such that
\[
S \cap U = \{ x \in U : \phi_{s+1}(x) = 0, \ldots, \phi_{r-1}(x) = 0 \}.
\]
By replacing $U$ by the connected component of $U$ containing $p$ we
shall henceforth assume that $U$ is connected. Let $T$ be the
$s$-dimensional linear subspace
\[
T = \{ y = (y_1, \ldots, y_{r-1}) \in \R^{r-1} :
y_{s+1} = 0, \ldots, y_{r-1} = 0 \}
\]
of $\R^{r-1}$. Then the image of $S \cap U$ under $\Phi$ is
$T \cap V$. Let
\[
f(x, x_r) = f_0(x) x_r^d + f_1(x) x_r^{d-1} + \cdots + f_d(x),
\]
where $f_0(x) \ne 0$. We shall denote by $g(y, y_r)$ the transform of
$f(x, x_r)$ by $\Phi$. That is, if $y \in V$, then
$g_j(y) = f_j(\Phi^{-1}(y))$ for $0 \le j \le d$, and
\[
g(y, y_r) = g_0(y) y_r^d + g_1(y) y_r^{d-1} + \cdots + g_d(y).
\]
As the analytic map $\Phi^{-1}$ is not in general a polynomial map,
$g(y, y_r)$ is not in general a polynomial in $y$ and $y_r$. However
we can regard $g(y, y_r)$ as a polynomial in $y_r$ whose coefficients
$g_j(y)$ are analytic functions of $y$, defined in $V$ (so
$g(y, y_r)$ is an instance of the real analogue of the
pseudopolynomials introduced in Sec.~2.1). Note that $y_r = 0$ is a
root of multiplicity $m$ of $g(0, y_r)$, that $g(0, \alpha) \ne 0$ for
complex $\alpha$ satisfying $0 < |\alpha| \le \epsilon$, and that,
for each fixed $y \in V$, $g(y, y_r)$ has exactly $m$ roots
(multiplicities counted) inside $C$. Let $E(y)$ be the discriminant
of $g(y, y_r)$. Then $E(y)$ is an analytic function in $V$ and
$E(y) = D(\Phi^{-1}(y))$ for every $y \in V$. Recall that $D$ is
order-invariant in $S$. Hence, by the remark following the definition
of order-invariant in Section~3.2, $E$ is order-invariant in
$T \cap V$.

We should like to find a box $B = B(0; \delta) \subseteq V$ about $0$
in $\R^{r-1}$, where $\delta = (\delta_1, \ldots, \delta_{r-1})$, and
an analytic function $\eta(y_1, \ldots, y_s)$ from
$B^{(s)} = \{ (y_1, \ldots, y_s) \in \R^s : |y_i| < \delta_i,\
1 \le i \le s \}$ into $(-\epsilon, +\epsilon)$, such that for every
$y = (y_1, \ldots, y_s, 0, \ldots, 0) \in T \cap B$ and every
$\alpha \in (-\epsilon, +\epsilon)$,
\begin{equation}
g(y, \alpha) = 0 \text{ if and only if } \alpha = \eta(y_1, \ldots, y_s) \tag{3.3.4}
\end{equation}
and such that $g$ is order-invariant in the set
\begin{equation}
\begin{aligned}
G = \{ (y, y_r) \in \R^r :\ &y = (y_1, \ldots, y_s, 0, \ldots, 0) \in T \cap B \\
&\text{and } y_r = \eta(y_1, \ldots, y_s) \}.
\end{aligned}
\tag{3.3.5}
\end{equation}
(Note that as $B^{(s)} \equiv T \cap B$ we can think of $\eta$ as a map
from $T \cap B$ into $(-\epsilon, +\epsilon)$, and the set $G$ as the
graph of $\eta$.) We shall make an excursion into the complex domain
from which we shall return with the box $B$ and the function $\eta$.
Upon detransforming $B$ and $\eta$ via $\Phi$ we shall have our
desired $N$ and $\theta$.

Choose $\rho = (\rho_1, \ldots, \rho_{r-1}) \in V$, with $\rho_i > 0$
for all $i$, such that $B_1 := B(0; \rho) \subseteq V$ and the power
series expansion about $0$ of every $g_j$ is absolutely convergent at
$y = \rho$. By Theorem~2.1.2, the power series expansion about $0$ of
every $g_j$, considered as a complex power series in the complex
variables $z_1, \ldots, z_{r-1}$, is absolutely convergent in the
polydisc $\Delta_1 := \Delta(0; \rho) \subseteq \C^{r-1}$ and the sum
of the series, $g_j(z) = g_j(z_1, \ldots, z_{r-1})$, is holomorphic in
$\Delta_1$. For $z \in \Delta_1$, set
\[
g(z, z_r) = g_0(z) z_r^d + g_1(z) z_r^{d-1} + \cdots + g_d(z)
\]
(so $g(z, z_r)$ is a pseudopolynomial in $\Delta_1$) and let $E(z)$
be the discriminant of $g(z, z_r)$. Let
\begin{equation}
E(y) = \sum_{i_1, \ldots, i_{r-1} \ge 0}
e_{i_1, \ldots, i_{r-1}} y_1^{i_1} \cdots y_{r-1}^{i_{r-1}} \tag{3.3.6}
\end{equation}
be the power series expansion for $E(y)$ about $0$. Then the power
series expansion for $E(z)$ about $0$, which is absolutely convergent
in $\Delta_1$, is obtained by substituting $z_1, \ldots, z_{r-1}$ for
$y_1, \ldots, y_{r-1}$ in (3.3.6).

Let $T^*$ be the $s$-dimensional linear subspace
\[
\{ z = (z_1, \ldots, z_{r-1}) \in \C^{r-1} :
z_{s+1} = 0, \ldots, z_{r-1} = 0 \}
\]
of $\C^{r-1}$ (note that $s$ is the dimension of $T^*$ considered as a
complex subspace of the complex vector space $\C^{r-1}$). We claim
that, after refining $\Delta_1$ suitably, we can assume that $E(z)$ is
order-invariant in $T^* \cap \Delta_1$. We prove this claim now. Let
$\ord_0 E = \mu$. Then every partial derivative of $E(z)$ of order
less than $\mu$ vanishes at $0$. Let
\[
P(z) = \frac{\partial^{i_1 + \cdots + i_{r-1}}}
{\partial z_1^{i_1} \cdots \partial z_{r-1}^{i_{r-1}}} E(z)
\]
be a partial derivative of $E(z)$ of order
$\sum_{j=1}^{r-1} i_j < \mu$. By Theorem~2.1.2, $P(z)$ is holomorphic
in $\Delta_1$ and its power series expansion about $0$, which is
absolutely convergent in $\Delta_1$, is obtained by differentiating
(3.3.6) (with the $y_i$'s replaced by $z_i$'s) term-by-term. Let
\[
Q(y) = \frac{\partial^{i_1 + \cdots + i_{r-1}}}
{\partial y_1^{i_1} \cdots \partial y_{r-1}^{i_{r-1}}} E(y).
\]
By Theorem~2.1.2, $Q(y)$ is analytic in $B_1$, and its power series
expansion about $0$, which is absolutely convergent in $B_1$, is
obtained by differentiating (3.3.6) term-by-term. Hence the power
series for $P(z)$ can be obtained by substituting
$z_1, \ldots, z_{r-1}$ for $y_1, \ldots, y_{r-1}$ in the power series
for $Q(y)$. Recall that $E(y)$ is order-invariant in $T \cap V$.
Hence, as $\ord_0 E = \mu$, $Q(y)$ vanishes throughout $T \cap B_1$.
It follows that substituting $y_{s+1} = 0, \ldots, y_{r-1} = 0$ in
the power series for $Q(y)$ yields the zero power series in
$y_1, \ldots, y_s$. Hence substituting $z_{s+1} = 0, \ldots,
z_{r-1} = 0$ in the power series for $P(z)$ yields the zero power
series in $z_1, \ldots, z_s$. It follows that $P(z)$ vanishes
throughout $T^* \cap \Delta_1$. We have shown that every partial
derivative of $E(z)$ of order less than $\mu$ vanishes throughout
$T^* \cap \Delta_1$. As $\ord_0 E = \mu$, some partial derivative of
$E(z)$ of order $\mu$, say $R(z)$, does not vanish at $0$. Let
$\Delta_2 \subseteq \Delta_1$ be a polydisc about $0$ in which
$R \ne 0$. Then $\ord_z E = \mu$ for every $z \in T^* \cap \Delta_2$,
so $E(z)$ is order-invariant in $T^* \cap \Delta_2$. Replacing
$\Delta_1$ by $\Delta_2$ we may assume that $E(z)$ is order-invariant
in $T^* \cap \Delta_1$.

We wish to study the structure of the zero set of $g(z, z_r)$ near
the origin. An application of the Weierstrass preparation theorem will
facilitate this study. Recall that $g(z, z_r)$ is holomorphic in the
polydisc $\Delta_1 \times \Delta(0; \epsilon)$, that $z_r = 0$ is a
root of $g(0, z_r)$ of multiplicity $m$, and that $g(0, z_r) \ne 0$
for $0 < |z_r| \le \epsilon$. By the Weierstrass preparation theorem
(Theorem~2.1.10), there is a polydisc
$\Delta_2 \subseteq \Delta_1$, a function $u(z, z_r)$ holomorphic and
non-vanishing in $\Delta' := \Delta_2 \times \Delta(0; \epsilon)$,
and a Weierstrass polynomial
\[
h(z, z_r) = z_r^m + a_1(z) z_r^{m-1} + \cdots + a_m(z)
\]
in $\Delta_2$, such that
\begin{equation}
g(z, z_r) = u(z, z_r) h(z, z_r) \tag{3.3.7}
\end{equation}
for all $(z, z_r) \in \Delta'$, and such that for each fixed
$z \in \Delta_2$, all the $m$ roots of $h(z, z_r)$ are contained in
the disc $\Delta(0; \epsilon)$. By (3.3.7), as $u \ne 0$ in $\Delta'$,
the zero set of $g(z, z_r)$ is the same as that of $h(z, z_r)$, in
$\Delta'$.

We claim that $u(z, z_r)$ is in fact a pseudopolynomial in $\Delta_2$.
We prove this claim now. Let $z$ be a fixed point of $\Delta_2$. Let
$z_r = \xi$ be a root of $h(z, z_r)$. Then $\xi \in \Delta(0; \epsilon)$.
Hence, as (3.3.7) holds near $z_r = \xi$ in the $z_r$-plane, and as
$u(z, z_r) \ne 0$ near $z_r = \xi$, $g(z, \xi) = 0$ and the
multiplicity of the zero $\xi$ of $h(z, z_r)$ is equal to the
multiplicity of the zero $\xi$ of $g(z, z_r)$. We have shown that
every root of $h(z, z_r)$ is a root of $g(z, z_r)$, and moreover has
the same multiplicity in $g(z, z_r)$ as it has in $h(z, z_r)$. Thus
$h(z, z_r) \mid g(z, z_r)$ in $\C[z_r]$ and the quotient of
$g(z, z_r)$ by $h(z, z_r)$ is clearly equal to $u(z, z_r)$. Let $d$
be the degree of $g(z, z_r)$. Then $m \le d$, and the degree of
$u(z, z_r)$ is $d - m$. Let us write
\[
u(z, z_r) = u_0(z) z_r^d + u_1(z) z_r^{d-1} + \cdots + u_d(z),
\]
where $u_0(z) = \cdots = u_{m-1}(z) = 0$, and $u_{d-m}(z) \ne 0$.
Letting $z$ now be a variable point of $\Delta_2$, the coefficients
$u_j(z)$ of $u(z, z_r)$ can be regarded as functions defined in
$\Delta_2$. That these functions are in fact holomorphic in $\Delta_2$
is seen by equating coefficients in (3.3.7). We have established our
claim that $u(z, z_r)$ is a pseudopolynomial in $\Delta_2$.

Let $F(z)$ be the discriminant of $h(z, z_r)$. We shall find a
function $Q(z)$ holomorphic in $\Delta_2$ such that
\begin{equation}
E(z) = Q(z) F(z) \tag{3.3.8}
\end{equation}
for all $z$ in $\Delta_2$. If $d > m$ in which case $u(z, z_r)$ has
positive degree in $z_r$, then set
\[
Q(z) = G(z) R(z)^2,
\]
where $G(z)$ is the discriminant of $u(z, z_r)$ and $R(z)$ is the
resultant of $u(z, z_r)$ and $h(z, z_r)$. Equation (3.3.8) holds by
Theorem~2.3.3. If $d = m$ in which case $u(z, z_r)$ has degree zero
in $z_r$, then set
\[
Q(z) = (u_0(z))^{2m - 2}.
\]
Equation (3.3.8) holds by Equation (2.3.6).

By Lemma~3.2.2, $F(z)$ is order-invariant in $T^* \cap \Delta_2$.
Hence, by Zariski's 1975 theorem (Theorem~4.1.1), there exists a
polydisc $\Delta_3 \subseteq \Delta_2$ about $0$, and a holomorphic
function $\psi(z_1, \ldots, z_s)$ from $\Delta_3^{(s)}$ into
$\Delta(0; \epsilon)$ such that for all
$z = (z_1, \ldots, z_s, 0, \ldots, 0) \in T^* \cap \Delta_3$ and all
$\alpha \in \Delta(0; \epsilon)$,
\begin{equation}
h(z, \alpha) = 0 \text{ if and only if } \alpha = \psi(z_1, \ldots, z_s) \tag{3.3.9}
\end{equation}
and such that $h$ is order-invariant in the set
\[
\begin{aligned}
G^* = \{ (z, z_r) \in \C^r :\
&z = (z_1, \ldots, z_s, 0, \ldots, 0) \in T^* \cap \Delta_3 \\
&\text{and } z_r = \psi(z_1, \ldots, z_s) \}.
\end{aligned}
\]
As remarked following the statement of Zariski's theorem in Sec.~4.1,
$\Delta_3^{(s)}$ can be identified with $T^* \cap \Delta_3$, and
$\psi$ regarded as a map from $T^* \cap \Delta_3$ into
$\Delta(0; \epsilon)$, with graph $G^*$.

Let the power series expansion for $\psi(z_1, \ldots, z_s)$ about $0$
be absolutely convergent at the point
$(z_1, \ldots, z_s) = (\beta_1, \ldots, \beta_s)$ of
$\Delta_3^{(s)}$, where $\beta_i > 0$ for each $i$. Choose
$\beta_{s+1}, \ldots, \beta_{r-1} > 0$ so that, if
$\delta = (\beta_1, \ldots, \beta_s, \beta_{s+1}, \ldots, \beta_{r-1})$,
then the polydisc $\Delta_4 := \Delta(0; \delta)$ in $\C^{r-1}$ is
contained in $\Delta_3$. Let $B = B(0; \delta)$ in $\R^{r-1}$ and let
$\eta$ be the restriction of $\psi$ to $B^{(s)}$. Then $B \subseteq V$
and we claim that $\eta$ is a real-valued analytic function, that for
every $y = (y_1, \ldots, y_s, 0, \ldots, 0) \in T \cap B$ and every
$\alpha \in (-\epsilon, +\epsilon)$, (3.3.4) holds, and that
$g(y, y_r)$ is order-invariant in ``the graph of $\eta$'' (i.e.,
strictly speaking, in the set $G$ defined by (3.3.5)). We prove these
claims now. Let $(y_1, \ldots, y_s) \in B^{(s)}$ and let
$y = (y_1, \ldots, y_s, 0, \ldots, 0)$ be the corresponding element of
$T \cap B$. By (3.3.9), $\alpha := \psi(y_1, \ldots, y_s)$ is a root
of $h(y, y_r)$ inside $\Delta(0; \epsilon)$. Hence, by (3.3.7),
$\alpha$ is a root of $g(y, y_r)$ inside $\Delta(0; \epsilon)$. But
$g(y, y_r)$ is a real polynomial in $y_r$. Hence $\bar{\alpha}$ is a
root of $g(y, y_r)$ inside $\Delta(0; \epsilon)$, and so by (3.3.7),
$h(y, \bar{\alpha}) = 0$. But (3.3.9) implies that $\bar{\alpha} =
\alpha$, and hence $\alpha$ is a real number. Thus $\eta$ is a
real-valued function, and $\eta : B^{(s)} \to (-\epsilon, +\epsilon)$.
Let
\begin{equation}
\psi(z_1, \ldots, z_s) = \sum_{j \ge 0} d_j z_1^{j_1} \cdots z_s^{j_s}, \tag{3.3.10}
\end{equation}
where $j$ denotes the multi-index $(j_1, \ldots, j_s)$, be the power
series expansion for $\psi(z_1, \ldots, z_s)$ about $0$. Then
(3.3.10) is absolutely convergent in $\Delta_4^{(s)}$, by
Theorem~2.1.2. Let $d_j = b_j + i c_j$, where $b_j$ and $c_j$ are
real numbers, for each multi-index $j$. By Theorem 48 of Chap.~6,
[KAP52],
\[
\psi(\beta_1, \ldots, \beta_s) = \sum_{j \ge 0} b_j \beta_1^{j_1} \cdots \beta_s^{j_s}
+ i \left( \sum_{j \ge 0} c_j \beta_1^{j_1} \cdots \beta_s^{j_s} \right),
\]
where both the real and imaginary parts of the right-hand side are
absolutely convergent. By Theorem~2.1.2, each of the series
\[
\sum_{j \ge 0} b_j y_1^{j_1} \cdots y_s^{j_s}, \quad
\sum_{j \ge 0} c_j y_1^{j_1} \cdots y_s^{j_s}
\]
is absolutely convergent in $B^{(s)}$, and the sum of each series is
an analytic function in $B^{(s)}$. But for
$(y_1, \ldots, y_s) \in B^{(s)}$,
\[
\begin{aligned}
\psi(y_1, \ldots, y_s) = \eta(y_1, \ldots, y_s)
&= \sum_{j \ge 0} d_j y_1^{j_1} \cdots y_s^{j_s} \\
&= \sum_{j \ge 0} b_j y_1^{j_1} \cdots y_s^{j_s}
+ i \left( \sum_{j \ge 0} c_j y_1^{j_1} \cdots y_s^{j_s} \right) \\
&= \sum_{j \ge 0} b_j y_1^{j_1} \cdots y_s^{j_s}, \quad
\text{as $\eta$ is real-valued.}
\end{aligned}
\]
Hence $\eta$ is analytic in $B^{(s)}$. It is straightforward to verify
the remaining claims about $\eta$.

Finally, set $N = \Phi^{-1}(B)$ and define
$\theta : S \cap N \to (-\epsilon, +\epsilon)$ by
$\theta(x) = (\eta \circ \phi)(x)$, where $\phi : S \cap N \to B^{(s)}$
is the chart for $S$ corresponding to $\Phi$. As
$\theta \circ \phi^{-1} = \eta$ is analytic in $B^{(s)}$, $\theta$ is
analytic in $S \cap N$. Clearly (3.3.1) holds for every
$x \in S \cap N$ and every $\alpha \in (-\epsilon, +\epsilon)$, and
the order-invariance of $f$ in the graph of $\theta$ follows from that
of $g$ in the graph of $\eta$, by the remark following the definition
of order-invariance in Section~3.2. This completes the proof of
Assertion~2 and hence of the lifting theorem.
\end{proof}

\newpage

\section*{Chapter Four: The Zariski Theorem}

This chapter presents an exposition of Zariski's 1975 theorem on
equisingularity over the complex field [ZAR75]. Our presentation of
this result is self-contained and in many respects quite different
from Zariski's. Our basic tools are analytic functions of several
complex variables, and elementary topology (covering spaces in
particular). Puiseux series (or fractional power series) enter into
one aspect of the argument.

The proof is divided into two cases: the codimension one case, and the
remaining case. The codimension one case was established by Zariski in
1965 [ZAR65] (the essential ideas were known earlier [ZAR35]).
Section~2 contains an exposition of this result. The remaining case
is proved by reduction to the codimension one case. This reduction is
the essential content of [ZAR75]. Section~1 provides our version of
the reduction.

Throughout this chapter we shall denote the $(n-1)$-tuple
$(z_1, \ldots, z_{n-1})$ by $z$.

\section{The General Theorem}

The following theorem is an adaptation of a result published by
Zariski in 1975 [ZAR75]:

\begin{theorem*}[4.1.1]
Let $h(z, z_n)$ be a Weierstrass polynomial of positive degree in the
polydisc $\Delta_1$ about $0$ in $\C^{n-1}$, and assume that for every
fixed $z$ in $\Delta_1$, all the roots of $h(z, z_n)$ are contained in
the disc $\Delta(0; \epsilon)$ about $0$ in $\C$. Let $F(z)$ be the
discriminant of $h(z, z_n)$ and assume that $F(z)$ does not vanish
identically. Let $1 \le s \le n-2$, let
\[
T^* = \{ z = (z_1, \ldots, z_{n-1}) \in \C^{n-1} :
z_{s+1} = 0, \ldots, z_{n-1} = 0 \},
\]
and assume that $F$ is order-invariant in $T^* \cap \Delta_1$. Then
there exists a polydisc $\Delta_2 \subseteq \Delta_1$ about $0$, and
a holomorphic function $\psi(z_1, \ldots, z_s)$ from $\Delta_2^{(s)}$
into $\Delta(0; \epsilon)$, such that for all
$z = (z_1, \ldots, z_s, 0, \ldots, 0)$ in $T^* \cap \Delta_2$ and all
$\alpha$ in $\Delta(0; \epsilon)$,
\[
h(z, \alpha) = 0 \text{ if and only if } \alpha = \psi(z_1, \ldots, z_s),
\]
and such that $h$ is order-invariant in the set
\[
\begin{aligned}
G^* = \{ (z, z_n) \in \C^{n-1} \times \C :\
&z = (z_1, \ldots, z_s, 0, \ldots, 0) \in T^* \cap \Delta_2 \\
&\text{and } z_n = \psi(z_1, \ldots, z_s) \}.
\end{aligned}
\]
\end{theorem*}

\textit{Remark 1.} As $\Delta_2^{(s)}$ can be identified with
$T^* \cap \Delta_2$ under the mapping
$\iota(z_1, \ldots, z_s) = (z_1, \ldots, z_s, 0, \ldots, 0)$, we can
think of $\psi$ as a map from $T^* \cap \Delta_2$ into
$\Delta(0; \epsilon)$, and the set $G^*$ as the graph of $\psi$.

\textit{Remark 2.} Zariski's theorem itself actually allows in place
of the particular $s$-dimensional linear subspace $T^*$ of $\C^{n-1}$
an arbitrary $s$-dimensional submanifold $S^*$ of $\C^{n-1}$
(containing $0$). We shall not need to formulate this slightly more
general theorem here, however.

\begin{proof}[Proof of 4.1.1]
Our proof of the theorem is given in numbered stages.

\begin{enumerate}
\item Let $\ord_0 F = r$. We may assume without loss of generality
that $\dfrac{\partial^r F}{\partial z_{n-1}^r}(0) \ne 0$. (If this is
not the case then we can find an invertible linear transformation $L$
of $\C^{n-1}$ fixing $T^*$ pointwise such that if $F' := F \circ L^{-1}$
is the transform of $F$ by $L$ then
$\dfrac{\partial^r F'}{\partial z_{n-1}^r}(0) \ne 0$.)

\item For the remainder of this section we shall denote the $s$-tuple
of variables $(y_1, \ldots, y_s)$ by $y$. Let the power series
expansion for $F(z)$ about $0$ be
\begin{equation}
\sum_{i_1, \ldots, i_{n-1} \ge 0} f_{i_1, \ldots, i_{n-1}}
z_1^{i_1} \cdots z_{n-1}^{i_{n-1}}. \tag{4.1.1}
\end{equation}
Then $f_{i_1, \ldots, i_{n-1}} = 0$ for all
$(i_1, \ldots, i_{n-1})$ with $\sum_{j=1}^{n-1} i_j < r$, and
$f_{0, \ldots, 0, r} \ne 0$ (by (1)). The series is absolutely
convergent in $\Delta_1$.

\item The series (4.1.1) can be arranged as an iterated series thus:
\begin{equation}
\sum_{i_{s+1}, \ldots, i_{n-1} \ge 0}
\left( \sum_{i_1, \ldots, i_s \ge 0}
f_{i_1, \ldots, i_{n-1}} z_1^{i_1} \cdots z_s^{i_s} \right)
z_{s+1}^{i_{s+1}} \cdots z_{n-1}^{i_{n-1}}. \tag{4.1.2}
\end{equation}
For each fixed $i_{s+1}, \ldots, i_{n-1} \ge 0$, the inner series in
(4.1.2) is absolutely convergent in $\Delta_1^{(s)}$. Let
$p_{i_{s+1}, \ldots, i_{n-1}}(z)$ denote the sum of this inner series.
Let $\delta = (\delta_1, \ldots, \delta_{n-1})$ be the polyradius of
$\Delta_1$. Then the outer series
\begin{equation}
\sum_{i_{s+1}, \ldots, i_{n-1} \ge 0}
p_{i_{s+1}, \ldots, i_{n-1}}(z)
z_{s+1}^{i_{s+1}} \cdots z_{n-1}^{i_{n-1}} \tag{4.1.3}
\end{equation}
is absolutely convergent in the polydisc
$|z_{s+1}| < \delta_{s+1}, \ldots, |z_{n-1}| < \delta_{n-1}$.

\item For $z \in \Delta_1^{(s)}$ and $k \ge 0$, let
\[
P_k(z, z_{s+1}, \ldots, z_{n-1}) = \sum_{i_{s+1} + \cdots + i_{n-1} = k}
p_{i_{s+1}, \ldots, i_{n-1}}(z)
z_{s+1}^{i_{s+1}} \cdots z_{n-1}^{i_{n-1}}.
\]
Then $P_0 = \cdots = P_{r-1} = 0$, by order-invariance of $F$ in
$T^* \cap \Delta_1$. Hence, for each fixed $z \in \Delta_1^{(s)}$, and
for $|z_{s+1}| < \delta_{s+1}, \ldots, |z_{n-1}| < \delta_{n-1}$,
\[
F(z, z_{s+1}, \ldots, z_{n-1}) = P_r(z, z_{s+1}, \ldots, z_{n-1})
+ P_{r+1}(z, z_{s+1}, \ldots, z_{n-1}) + \cdots
\]
(grouping terms in the series (4.1.3)).

\item \textbf{Case I: $s = n - 2$.}

We have
\[
F(z, z_{n-1}) = p_r(z) z_{n-1}^r + p_{r+1}(z) z_{n-1}^{r+1} + \cdots.
\]
Hence $F(z, z_{n-1}) = z_{n-1}^r N(z, z_{n-1})$, where
\[
N(z, z_{n-1}) = p_r(z) + p_{r+1}(z) z_{n-1} + \cdots,
\]
for $(z, z_{n-1}) \in \Delta_1$. It can be shown that $N$ is
holomorphic. Now $N(0, 0) = p_r(0) = f_{0, \ldots, 0, r} \ne 0$ (from
(2)). Hence, there exists $\Delta_2 \subseteq \Delta_1$ such that
$N \ne 0$ in $\Delta_2$. We can therefore apply Theorem~4.2.1 to
obtain the required function $\psi$.

\item \textbf{Case II: $1 \le s \le n - 3$.}

Our essential strategy is to reduce Case II to Case I. We use a
quadratic transformation (or ``blowing-up map'') to bring about the
reduction to the codimension~1 case. Define the map
$Q : \C^{n-1} \to \C^{n-1}$ by
\[
Q(z, z_{s+1}, \ldots, z_{n-1}) = (z, z_{s+1} z_{n-1}, \ldots,
z_{n-2} z_{n-1}, z_{n-1}).
\]
$Q$ is called a quadratic transformation, as each of its components
has degree at most two, and some component has degree exactly two.
$Q$ maps the entire hyperplane $H^* : z_{n-1} = 0$ to the
$s$-dimensional linear subspace $T^*$ of $\C^{n-1}$. Even though
strictly speaking $Q^{-1}$ does not exist, one can think of a ``map''
$Q^{-1}$ which sends a point
$z = (z, z_{s+1}, \ldots, z_{n-1})$, with $z_{n-1} \ne 0$, to the
point
\[
Z = \left( z, \frac{z_{s+1}}{z_{n-1}}, \ldots,
\frac{z_{n-2}}{z_{n-1}}, z_{n-1} \right)
\]
(for which $Q(Z) = z$), and which sends a point $z \in T^*$ to the
\emph{set} of points $Z \in \C^{n-1}$ for which $Q(Z) = z$. Points in
$\{ z_{n-1} = 0 \} - T^*$ would have no image under $Q^{-1}$, however.
We say that $T^*$ is \emph{blown-up} by $Q^{-1}$ to $H^*$, and call
$Q^{-1}$ a \emph{blowing-up map}.

Where $\delta_i' = \delta_i$ for $1 \le i \le s$ and $i = n-1$, and
$\delta_i' = \delta_i / \delta_{n-1}$ for $s+1 \le i \le n-2$, and
$\delta' = (\delta_1', \ldots, \delta_{n-1}')$, let
$\Delta_1' = \Delta(0; \delta')$. We have that
$Q(\Delta_1') \subseteq \Delta_1$. Hence if we set
$h'(Z, Z_n) = h(Q(Z), Z_n)$ for $Z \in \Delta_1'$, then $h'$ is a
Weierstrass polynomial in $\Delta_1'$. Let $F'(Z)$ be the discriminant
of $h'(Z, Z_n)$.

\item \textbf{Case II (continued).}

We show that $F'$ satisfies the hypotheses of Theorem~4.2.1 (and
equivalently that $F'$ is order-invariant in $H^*$, near $0$). Let
$Z = (z, Z_{s+1}, \ldots, Z_{n-1})$ be an element of $\Delta_1'$, and
let
\[
N(Z) = P_r(z, Z_{s+1}, \ldots, Z_{n-2}, 1)
+ Z_{n-1} P_{r+1}(z, Z_{s+1}, \ldots, Z_{n-2}, 1) + \cdots.
\]
Then
\[
\begin{aligned}
F'(Z) = F(Q(Z)) &= F(z, Z_{s+1} Z_{n-1}, \ldots, Z_{n-2} Z_{n-1}, Z_{n-1}) \\
&= P_r(z, Z_{s+1} Z_{n-1}, \ldots, Z_{n-2} Z_{n-1}, Z_{n-1}) \\
&\quad + P_{r+1}(z, Z_{s+1} Z_{n-1}, \ldots, Z_{n-2} Z_{n-1}, Z_{n-1}) + \cdots \\
&= Z_{n-1}^r N(Z),
\end{aligned}
\]
as
\[
P_k(z, Z_{s+1} Z_{n-1}, \ldots, Z_{n-2} Z_{n-1}, Z_{n-1}) =
Z_{n-1}^k P_k(z, Z_{s+1}, \ldots, Z_{n-1}, 1).
\]
It can be shown that $N$ is holomorphic; further, we have
\[
N(0) = P_r(0, 0, \ldots, 0, 1) = p_{0, \ldots, 0, r}(0) =
f_{0, \ldots, 0, r} \ne 0,
\]
by (2). Hence there exists $\Delta_2' \subseteq \Delta_1'$ such that
$N \ne 0$ in $\Delta_2'$.

\item \textbf{Case II (conclusion).}

By Theorem~4.2.1, there exists a polydisc $\Delta_3'$ contained in
$\Delta_2'$ and a holomorphic function $\psi'(Z_1, \ldots, Z_{n-2})$
from $\Delta_3'^{(n-2)}$ into $\Delta(0; \epsilon)$ such that for all
$Z = (Z_1, \ldots, Z_{n-2}, 0)$ in $H^* \cap \Delta_3'$ and all $Z_n$
in $\Delta(0; \epsilon)$,
\[
h'(Z, Z_n) = 0 \text{ if and only if }
Z_n = \psi'(Z_1, \ldots, Z_{n-2})
\]
and such that $h'$ is order-invariant in the set
\[
\begin{aligned}
K^* = \{ (Z, Z_n) \in \C^n :\
&Z = (Z_1, \ldots, Z_{n-2}, 0) \in H^* \cap \Delta_3' \\
&\text{and } Z_n = \psi'(Z_1, \ldots, Z_{n-2}) \}.
\end{aligned}
\]
Let $\nu'$ be the polyradius of $\Delta_3'$, let $\nu = Q(\nu')$,
and let $\Delta_3 = \Delta(0; \nu)$. Define
$\psi : \Delta_3^{(s)} \to \Delta(0; \epsilon)$ by
$\psi(z) = \psi'(z, 0, \ldots, 0)$. We shall show that $\psi$
satisfies the properties asserted of it in the conclusion of
Theorem~4.1.1.

It is clearly the case that for all $z = (z, 0, \ldots, 0)$ in
$T^* \cap \Delta_3$ and all $z_n$ in $\Delta(0; \epsilon)$,
\[
h(z, z_n) = 0 \text{ if and only if } z_n = \psi(z).
\]
It remains to show that $h$ is order-invariant in the set $G^*$
defined in the statement of the theorem (with `$\Delta_2$' replaced
by `$\Delta_3$').

Let $(w, w_n) \in G^*$ and let $t = \ord_{(w, w_n)} h$. We shall prove
that $t \ge \ord_{(w, w_n)} h'$. The order-invariance of $h$ in $G^*$
will then follow by the order-invariance of $h'$ in $K^*$. We have
$t \ge \ord_{(w, w_n)} h'$, by Theorem~2.3.1. We shall prove that
there exists a point $(w'', w_n) \in K^*$, with
$w'' = (w, w_{s+1}'', \ldots, w_{n-2}'', 0)$ an element of
$H^* \cap \Delta_3'$, such that $t \ge \ord_{(w'', w_n)} h'$. This
will show, by order-invariance of $h'$ in $K^*$, that
$t \ge \ord_{(w, w_n)} h'$.

Now $w$ is an element of $T^* \cap \Delta_3$, and so
$w = (\tilde{w}, 0, \ldots, 0)$, where $\tilde{w} \in \Delta_3^{(s)}$.
As $\ord_{(w, w_n)} h = t$, there is a term
\[
c (z_1 - w_1)^{e_1} \cdots (z_s - w_s)^{e_s}
z_{s+1}^{e_{s+1}} \cdots z_{n-1}^{e_{n-1}} (z_n - w_n)^{e_n},
\]
with $\sum_{j=1}^{n} e_j = t$ and $c \ne 0$, in the power series
expansion of $h$ about $(w, w_n)$. This term gives rise to the
non-zero term
\[
c (Z_1 - w_1)^{e_1} \cdots (Z_s - w_s)^{e_s}
Z_{s+1}^{e_{s+1}} \cdots Z_{n-2}^{e_{n-2}} Z_{n-1}^m
(Z_n - w_n)^{e_n}
\]
where $m = \sum_{j=s+1}^{n-1} e_j$, in the power series expansion of
$h'$ about $(w, w_n)$ (since $h'(Z, Z_n) = h(Q(Z), Z_n)$).

Let the power series expansion of $h'$ about $(w, w_n)$ be arranged
as an iterated series thus:
\[
\sum_{i_{s+1}, \ldots, i_n \ge 0}
q_{i_{s+1}, \ldots, i_n}(Z_{s+1}, \ldots, Z_{n-2})
(Z_1 - w_1)^{i_1} \cdots (Z_s - w_s)^{i_s}
Z_{n-1}^{i_{n-1}} (Z_n - w_n)^{i_n}
\]
(cf.\ stage~3 of the proof of Theorem~4.1.1). Now
$q_{e, m, e_n}$ is not identically zero, as it contains the non-zero
term $c Z_{s+1}^{e_{s+1}} \cdots Z_{n-2}^{e_{n-2}}$. Hence, there
exist $w_{s+1}'', \ldots, w_{n-2}''$, with
$w'' = (w, w_{s+1}'', \ldots, w_{n-2}'', 0)$ an element of
$H^* \cap \Delta_3'$, such that
\[
q := q_{e, m, e_n}(w_{s+1}'', \ldots, w_{n-2}'')
\]
is non-zero.

Now the power series expansion of $h'$ about $(w'', w_n)$ contains
the term
\[
q (Z_1 - w_1)^{e_1} \cdots (Z_s - w_s)^{e_s}
(Z_{s+1} - w_{s+1}'')^{e_{s+1}} \cdots (Z_{n-2} - w_{n-2}'')^{e_{n-2}}
Z_{n-1}^m (Z_n - w_n)^{e_n}.
\]
Hence, as $q \ne 0$ and the total degree of this term is
\[
e_1 + \cdots + e_s + m + e_n = t,
\]
we have $t \ge \ord_{(w'', w_n)} h'$. Hence, by the order-invariance
of $h'$ in $K^*$, we have $t \ge \ord_{(w, w_n)} h'$.

We have now shown that $t \ge \ord_{(w, w_n)} h'$. The
order-invariance of $h$ in $G^*$ now follows by the order-invariance
of $h'$ in $K^*$.
\end{enumerate}
\end{proof}

\section{Special Case (Codimension One)}

We present in this section a theorem established by Zariski in 1965
[ZAR65]. This theorem is essentially the special case $s = n-2$ of
Theorem~4.1.1. Our formulation of the result is adapted to our needs
and is somewhat different from Zariski's formulation. Our proof is
also different from the one Zariski gave. Throughout this section we
shall adopt the notational convention that $z = (z_1, \ldots, z_{n-2})$:
thus $z = (z, z_{n-1})$. Let $H^*$ be the hyperplane
$\{ z = (z, z_{n-1}) \in \C^{n-1} : z_{n-1} = 0 \}$ in $\C^{n-1}$. The
theorem can be stated as follows:

\begin{theorem*}[4.2.1 (Zariski)]
Let $n \ge 3$ and let $h(z, z_n)$ be a Weierstrass polynomial of
positive degree in the polydisc $\Delta_1$ about $0$ in $\C^{n-1}$,
and assume that for every fixed $z$ in $\Delta_1$, all the roots of
$h(z, z_n)$ are contained in the disc $\Delta(0; \epsilon)$ about $0$
in $\C$. Let $F(z)$ be the discriminant of $h(z, z_n)$ and suppose
that there exists a function $N(z)$, holomorphic and non-vanishing in
$\Delta_1$, and an integer $r \ge 0$, such that
\[
F(z) = z_{n-1}^r N(z)
\]
for all $z = (z, z_{n-1})$ in $\Delta_1$. Then there exists a polydisc
$\Delta_2 \subseteq \Delta_1$ about $0$, and a holomorphic function
$\psi(z)$ from $\Delta_2^{(n-2)}$ into $\Delta(0; \epsilon)$, such
that for every fixed $z = (z, 0)$ in $H^* \cap \Delta_2$ and every
$z_n$ in $\Delta(0; \epsilon)$,
\[
h(z, z_n) = 0 \text{ if and only if } z_n = \psi(z),
\]
and such that $h$ is order-invariant in the set
\begin{equation}
G^* = \{ (z, 0, z_n) \in \C^n : z \in \Delta_2^{(n-2)} \text{ and }
z_n = \psi(z) \}. \tag{4.2.1}
\end{equation}
\end{theorem*}

Note that the hypothesis on $F(z)$ that it be associated to the
function $z_{n-1}^r$ is essentially equivalent to the hypothesis that
$F(z)$ be order-invariant in $H^*$, near the origin. Thus the above
theorem is essentially the special case $s = n-2$ of Theorem~4.1.1:
that is, the case in which the linear subspace $T^*$ of $\C^{n-1}$
has codimension one.

This result of Zariski is equivalent to the following two theorems
taken together. The first theorem (4.2.2) asserts the ``nonsplitting''
of the locus $G^*$ of $h$ over $H^*$, and the second (4.2.3) asserts
the order-invariance of $h$ in $G^*$.

\begin{theorem*}[4.2.2]
Let $h(z, z_n)$ be a Weierstrass polynomial of positive degree in the
polydisc $\Delta_1$ about $0$ in $\C^{n-1}$, and assume that $h$
satisfies the hypotheses of Theorem~4.2.1. (In particular, assume
that for every fixed $z$ in $\Delta_1$, all the roots of $h(z, z_n)$
are contained in the disc $\Delta(0; \epsilon)$ in $\C$.) Then there
exists a holomorphic function $\psi(z)$ from $\Delta_1^{(n-2)}$ into
$\Delta(0; \epsilon)$ such that for every fixed
$z = (z, 0) \in H^* \cap \Delta_1$ and every $z_n \in \Delta(0; \epsilon)$,
\[
h(z, z_n) = 0 \text{ if and only if } z_n = \psi(z).
\]
\end{theorem*}

\begin{theorem*}[4.2.3]
Let $h(z, z_n)$ be a Weierstrass polynomial of degree $m \ge 1$ in the
polydisc $\Delta_1$ about $0$ in $\C^{n-1}$, and assume that $h$
satisfies the hypotheses of Theorem~4.2.1. Assume further that
$h(z, 0, z_n) = z_n^m$ for all $z \in \Delta_1^{(n-2)}$. Then there
exists a polydisc $\Delta_2 \subseteq \Delta_1$ about $0$ such that
$h$ is order-invariant in the set
$\Delta_2^{(n-2)} \times \{ (0, 0) \}$.
\end{theorem*}

\begin{proof}[Proof that 4.2.1 $\Leftrightarrow$ (4.2.2 \& 4.2.3)]
Clearly 4.2.1 implies 4.2.2 and 4.2.3. Assume conversely that 4.2.2
and 4.2.3 hold. Let $h(z, z_n)$ be a Weierstrass polynomial of degree
$m \ge 1$ in the polydisc $\Delta_1$ about $0$ in $\C^{n-1}$, and
assume that $h$ satisfies the hypotheses of Theorem~4.2.1.
Theorem~4.2.2 gives us the required function $\psi$, defined in
$\Delta_1^{(n-2)}$. It remains to establish that there is a polydisc
$\Delta_2 \subseteq \Delta_1$ about $0$ such that $h$ is
order-invariant in the set $G^*$ defined by equation~(4.2.1). For
$z = (z, z_{n-1}) \in \Delta_1$, let $h'(z, z_n) = h(z, z_n + \psi(z))$.
Then $h'(z, z_n)$ is a Weierstrass polynomial in $\Delta_1$ and
$h'(z, 0, z_n) = z_n^m$ for $z \in \Delta_1^{(n-2)}$. Moreover, by
root continuity, there exists a polydisc $\Delta_1' \subseteq \Delta_1$
about $0$ such that for every fixed $z \in \Delta_1'$, all the roots
of $h'(z, z_n)$ are contained in $\Delta(0; \epsilon)$. Hence, by
Theorem~4.2.3, there exists a polydisc $\Delta_2 \subseteq \Delta_1'$
about $0$ such that $h'$ is order-invariant in
$\Delta_2^{(n-2)} \times \{ (0, 0) \}$. Now the map
$(z, z_n) \mapsto (z, z_n + \psi(z))$ is an analytic isomorphism of
$\Delta_1 \times \C$ onto itself. Hence, by Theorem~2.3.1, the order
of $h'$ at $(z, 0, 0)$ is equal to the order of $h$ at
$(z, 0, \psi(z))$, for every $z \in \Delta_1^{(n-2)}$. Therefore $h$
is order-invariant in the set $G^*$ defined by equation~(4.2.1). We
have shown that Theorem~4.2.1 holds.
\end{proof}

It remains, therefore, to establish Theorems~4.2.2 and 4.2.3. We will
first need to study irreducible Weierstrass polynomials and
factorization. Let $R$ be an integral domain. A non-constant element
of the polynomial ring $R[z_n]$ is said to be \emph{reducible} if it
is a product of non-constant polynomials of lower degree, and is
otherwise said to be \emph{irreducible}. Let $R$ now be the ring of
all holomorphic functions in the polydisc $\Delta_1$ about $0$ in
$\C^{n-1}$. As $\Delta_1$ is connected, $R$ is an integral domain (by
the identity theorem, Theorem~2.1.4). Let $h(z, z_n)$ be a Weierstrass
polynomial in $\Delta_1$. Then $h$ can be regarded as a polynomial in
$z_n$ over the ring $R$, i.e.\ as an element of $R[z_n]$. It follows
by induction on the degree of $h$ that $h$ may be factored as a
product of irreducible Weierstrass polynomials thus:
\[
h = h_1 \cdots h_k.
\]
It will follow from Lemma~4.2.5 below that the $h_i$ are uniquely
determined, provided that the discriminant of $h$ does not vanish
identically.

Let $U$ be an open subset of $\C^n$. A continuous function
$\Gamma : [0, 1] \to U$, with $\Gamma(0) = w_0$ and $\Gamma(1) = w_1$,
is called a \emph{path} in $U$ from $w_0$ to $w_1$. Let $\Gamma$ and
$\Gamma'$ be paths in $U$, from $w_0$ to $w_1$, and from $w_1$ to
$w_2$, respectively. Then the composite path $\Gamma' * \Gamma$ from
$w_0$ to $w_2$ is defined as follows:
\[
(\Gamma' * \Gamma)(t) =
\begin{cases}
\Gamma(2t) & \text{if } 0 \le t \le 1/2 \\
\Gamma'(2t - 1) & \text{if } 1/2 \le t \le 1.
\end{cases}
\]
The \emph{reverse path} $\bar{\Gamma}$ of $\Gamma$ is defined by
$\bar{\Gamma}(t) = \Gamma(1 - t)$, $0 \le t \le 1$.

\begin{lemma*}[4.2.4]
Let $\Delta_1$ be a polydisc about $0$ in $\C^{n-1}$ and let
$h(z, z_n)$ be a Weierstrass polynomial of positive degree in
$\Delta_1$. Let $U$ be an open subset of $\Delta_1$ in which the
discriminant of $h$ does not vanish. Let $w$ and $w'$ be points of
$U$, not necessarily distinct, and let $\Gamma$ be a path in $U$
from $w$ to $w'$. Let $\alpha$ be a root of $h(w, z_n)$. Then there
is a unique path $\phi$ in $\C^1$ such that $\phi(0) = \alpha$ and
$h(\Gamma(t), \phi(t)) = 0$ for all $t \in [0, 1]$.
\end{lemma*}

\begin{proof}
Let the degree of $h$ be $m \ge 1$, let
$V = \{ (z, z_n) \in U \times \C : h(z, z_n) = 0 \}$, and let $w_0$ be
a point of $U$. By $m$ applications of the implicit function theorem,
there exists a neighbourhood $W \subseteq U$ of $w_0$ such that the
portion of $V$ contained in $W \times \C$ consists of the disjoint
graphs of $m$ holomorphic functions from $W$ into $\C$. Hence
([MUN75], Chapter~8) $V$ is an $m$-fold covering of $U$, with
covering map $\pi : V \to U$ given by $\pi(z, z_n) = z$. The existence
and uniqueness of $\phi$ then follows by the path lifting property
(Lemma~4.1 in Ch.~8 of [MUN75]).
\end{proof}

\textit{Remark.} With the notations of the above lemma, we set
$\Gamma_h[\alpha] := \phi(1)$ and say $\Gamma$ carries $\alpha$ into
$\phi(1)$ (via $h$).

The following is an important lemma about irreducible Weierstrass
polynomials:

\begin{lemma*}[4.2.5]
Let $\Delta_1$ be a polydisc about $0$ in $\C^{n-1}$ and let
$h(z, z_n)$ be an irreducible Weierstrass polynomial in $\Delta_1$.
Let $G(z)$ be the discriminant of $h(z, z_n)$, let $F(z)$ be a
holomorphic function in $\Delta_1$, not identically zero, such that
$G(z) = 0 \Rightarrow F(z) = 0$ for all $z$ in $\Delta_1$, and let
$U = \{ z \in \Delta_1 : F(z) \ne 0 \}$. Let $w \in U$ and let
$\alpha$ be a root of $h(w, z_n)$. Then for every $w' \in U$ and
every root $\alpha'$ of $h(w', z_n)$, there exists a path $\Gamma$ in
$U$ from $w$ to $w'$ such that $\Gamma_h[\alpha] = \alpha'$.
\end{lemma*}

\begin{proof}
A proof of this result, using slightly different notation and
terminology, is given in [BMA48], Ch.~9, Sec.~3, pp 194--198.
\end{proof}

\textit{Remark.} The above lemma implies that the irreducible factors
of a Weierstrass polynomial whose discriminant does not vanish
identically are uniquely determined.

We can now give the

\begin{proof}[Proof of Theorem 4.2.2]
Factor $h$ into a product of irreducible Weierstrass polynomials
$h_1 \cdots h_k$. Let $1 \le i \le k$, and let $G(z)$ be the
discriminant of $h_i(z, z_n)$. Recall that the discriminant $F(z)$ of
$h(z, z_n)$ has the form $F(z) = z_{n-1}^r N(z)$ for all
$z \in \Delta_1$, where $N(z)$ is holomorphic and non-vanishing in
$\Delta_1$. Now the zero set of $G$ is contained in that of $F$ by
Lemma~2.3.4, and is hence a subset of $H^* \cap \Delta_1$. But if
$G(w, 0) = 0$ for some $w \in \Delta_1^{(n-2)}$, then $G(z, 0) = 0$
for all $z \in \Delta_1^{(n-2)}$ by zero system continuity (see
[WHI72], Ch.~1, Lemma~5E). Hence the zero set of $G$ is either
empty, or the entire region $H^* \cap \Delta_1$.

We shall show that, for each point $w = (w, 0)$ in $H^* \cap \Delta_1$,
there exists exactly one distinct root of $h_i(w, z_n)$. This is
clearly true if $G \ne 0$ in $\Delta_1$, as in this case the degree of
$h_i$ is $1$. Suppose, on the other hand, that $G \equiv 0$ in
$H^* \cap \Delta_1$, and that, for some
$w = (w, 0) \in H^* \cap \Delta_1$, $h_i(w, z_n)$ has some $l \ge 2$
distinct roots, say $\alpha_1, \ldots, \alpha_l$, with multiplicities
$m_1, \ldots, m_l$ respectively. Choose disjoint open discs
$D_1, \ldots, D_l$ about $\alpha_1, \ldots, \alpha_l$, respectively.
By root continuity of the bivariate Weierstrass polynomial
$h_i(w, z_{n-1}, z_n)$, there exists a disc $D'$ about $0$ in the
plane of the $z_{n-1}$-coordinate such that
$\{w\} \times D' \subseteq \Delta_1$ and such that, for every fixed
$w_{n-1}$ in $D'$, there are exactly $m_j$ roots (counted according
to multiplicity) of $h_i(w, w_{n-1}, z_n)$ in $D_j$, $1 \le j \le l$.

Let $w_{n-1}'$ be a non-zero point of $D'$, let $w' = (w, w_{n-1}')$
and let $\alpha$ and $\beta$ be roots of $h(w', z_n)$ in $D_1$ and
$D_2$ respectively. Let $U = \{ z \in \Delta_1 : G(z) \ne 0 \}$.
Then $U = \Delta_1 - H^*$. By Lemma~4.2.5, there is a path
$\Gamma(t) = (\Gamma(t), \Gamma_{n-1}(t))$ in $U$ from $w'$ to
itself such that $\Gamma_{h_i}[\alpha] = \beta$. Let
$r = (r, r_{n-1})$ be the polyradius of $\Delta_1$. We can deform
$\Gamma$ in $U$ to a path $\Gamma'$ in the set defined by $z = w$,
$|z_{n-1}| < r_{n-1}$, by means of the path homotopy ([MUN75], Ch.~8)
\[
H(s, t) = ((1 - s) \Gamma(t) + s w, \Gamma_{n-1}(t))
\]
(for which $H(0, t) = \Gamma(t)$ and $H(1, t) = (w, \Gamma_{n-1}(t))$).
Furthermore, we can deform $\Gamma'$ in $U$ to a path $\Gamma''$
along the circle defined by $z = w$, $|z_{n-1}| = |w_{n-1}'|$ by
means of the path homotopy
\[
K(s, t) = \left( w, (1 - s) \Gamma_{n-1}(t)
+ s \frac{\Gamma_{n-1}(t) |w_{n-1}'|}{|\Gamma_{n-1}(t)|} \right).
\]
Let $\phi''$ be the path in $\C^1$ such that $\phi''(0) = \alpha$ and
$h_i(\Gamma''(t), \phi''(t)) = 0$ for all $t \in [0, 1]$, given by
Lemma~4.2.4. By Theorem~4.3 of Ch.~8 of [MUN75], $\phi''(1) = \beta$
(which implies $\Gamma''_{h_i}[\alpha] = \beta$). Yet, as $\phi''(t)$
is a root of $h_i(\Gamma''(t), z_n)$, we must have that
$\phi''(t) \in \bigcup_{j=1}^{l} D_j$, for all $t$. Hence, as
$\phi''$ is continuous, the $D_j$ are disjoint, and
$\phi''(0) \in D_1$, we must have that $\phi''(t) \in D_1$ for all
$t$. This contradicts $\phi''(1) = \beta$, as $\beta$ is an element
of $D_2$. Hence our assumption that $h_i(w, z_n)$ has more than one
distinct root must be false. Thus $h_i(w, z_n)$ has exactly one
distinct root, for each fixed $w$ in the set $H^* \cap \Delta_1$.

We can now show that $h(w, z_n)$ has exactly one distinct root, for
each fixed $w \in H^* \cap \Delta_1$. There is nothing further to
prove if $k = 1$, so assume $k > 1$. Let
$w \in H^* \cap \Delta_1$, let $1 \le i < j \le k$, and let
$\alpha_i$ and $\alpha_j$ be the roots of $h_i(w, z_n)$ and
$h_j(w, z_n)$ respectively, and assume that $\alpha_i \ne \alpha_j$.
Then the resultant $R(z)$ of $h_i(z, z_n)$ and $h_j(z, z_n)$ does
not vanish at $z = w$. However, the zero set of $R$ contains $0$, is
contained in $H^* \cap \Delta_1$ (as $R(z) = 0 \Rightarrow F(z) = 0$),
and is hence equal to the entire region $H^* \cap \Delta_1$ (by
zero system continuity), contradicting $R(w) \ne 0$. Therefore, we
must have $\alpha_i = \alpha_j$ after all. Hence $h(w, z_n)$ has
exactly one distinct root, say $\alpha$.

For $w \in \Delta_1^{(n-2)}$, let $\psi(w)$ denote the unique root
$\alpha$ of $h(w, 0, z_n)$. Then, where $m$ is the degree of $h$, we
have
\[
h(w, 0, z_n) = (z_n - \psi(w))^m.
\]
Thus, where $a_1(z)$ is the coefficient of $z_n^{m-1}$ in
$h(z, z_n)$, we have $a_1(w, 0) = -m \psi(w)$. Hence $\psi$ is
holomorphic in $\Delta_1^{(n-2)}$.
\end{proof}

We turn now to establishing Theorem~4.2.3. The theorem will follow
from three lemmas on parametrizing an irreducible Weierstrass
polynomial which satisfies the hypotheses of the theorem. The first
lemma asserts the existence of a parametrization for such a
Weierstrass polynomial.

\begin{lemma*}[4.2.6]
Let $h(z, z_n)$ be an irreducible Weierstrass polynomial of degree
$m \ge 1$ in the polydisc $\Delta_1 = \Delta(0; \delta)$ in
$\C^{n-1}$, and assume that $h$ satisfies the hypotheses of
Theorem~4.2.3. (In particular, assume that for every fixed $z$ in
$\Delta_1$, all the roots of $h(z, z_n)$ are contained in the disc
$\Delta(0; \epsilon)$ in $\C$.) Let $\rho = \delta$, let
$\rho_{n-1} = \delta_{n-1}^{1/m}$, let $\rho = (\rho, \rho_{n-1})$,
and let $\Delta_1' = \Delta(0; \rho)$. Then there is a holomorphic
function $\phi : \Delta_1' \to \Delta(0; \epsilon)$ such that for
every fixed $z = (z, z_{n-1}) \in \Delta_1$ and every $z_n$ in
$\Delta(0; \epsilon)$,
\[
h(z, z_n) = 0 \text{ if and only if there exists } u,\ |u| < \rho_{n-1},
\text{ such that}
\]
\[
\begin{cases}
z_{n-1} = u^m \\
z_n = \phi(z, u).
\end{cases}
\]
\end{lemma*}

\begin{proof}
Even though the hypotheses of Theorem~4.2.1 specify $n \ge 3$, the
hypotheses of the theorem, as indeed the present lemma, can be stated
for $n = 2$ as well. The case $n = 2$ of the present lemma is
Theorem~10A of [WHI72], Ch.~1. The hypotheses of Theorem~4.2.1
permit the extension of the result to higher dimensions.
\end{proof}

With the notation of the above lemma, we shall say that the pair of
equations
\[
\begin{cases}
z_{n-1} = u^m \\
z_n = \phi(z, u)
\end{cases}
\]
defines a \emph{parametrization} for $h$ (in $\Delta_1$). In such a
parametrization let $m_1$ denote the order of $\phi$ at $0$ in $u$.
Then we shall see (Lemma~4.2.8 below) that for all $z$ sufficiently
close to $0$,
\[
\ord_{(z, 0, 0)} h = \min(m, m_1).
\]
That $h$ is order-invariant in the set
$\Delta_2^{(n-2)} \times \{ (0, 0) \}$ (near the origin) will follow
from this. The next result will be used in proving Lemma~4.2.8.

Let $K$ be an algebraically closed field of characteristic zero. We
shall denote by $K\{x\}$ the field of all fractional power series in
the indeterminate $x$ with coefficients in $K$ (see [WAL50], Ch.~4,
Sec.~3, in which the notation $K(x)^*$ is used). That is, $K\{x\}$
consists of all elements $\theta$ of the form
\begin{equation}
\theta = a_h x^{h/m} + a_{h+1} x^{(h+1)/m} + \cdots \tag{4.2.2}
\end{equation}
where $m \ge 1$ and $h$ is an integer (possibly negative). If
$\theta$ is an element of $K\{x\}$ given by (4.2.2) and $a_h \ne 0$,
then we define the \emph{order} of $\theta$, denoted by $O(\theta)$,
to be $h/m$. Puiseux' theorem (Theorem~3.1 of [WAL50], Ch.~4) states
that the field $K\{x\}$ is algebraically closed. Thus every
nonconstant polynomial $g(x, y)$ over $K\{x\}$ can be factored into
linear factors over $K\{x\}$ thus:
\[
g(x, y) = \prod_{i=1}^{l} (y - \theta_i);
\]
$\{ \theta_1, \ldots, \theta_l \}$ is called the set of
\emph{Puiseux roots} of $g(x, y)$. We can now state and prove our
next result.

\begin{lemma*}[4.2.7]
Let $h(z, z_n)$ be an irreducible Weierstrass polynomial of degree
$m \ge 1$ in the polydisc $\Delta_1$ about $0$ (in $\C^{n-1}$), and
assume that $h$ satisfies the hypotheses of Theorem~4.2.3. Let
\[
\begin{cases}
z_{n-1} = u^m \\
z_n = \phi(z, u)
\end{cases}
\]
be a parametrization for $h$ and assume that $\phi$ is not
identically zero. Assume that the power series expansion for $\phi$
about $0$ is absolutely convergent in $\Delta_1'$ (using the notation
of the previous lemma) and let this power series be arranged as
follows:
\[
c_1(z) u^{m_1} + c_2(z) u^{m_2} + \cdots
\]
where $0 < m_1 < m_2 < \cdots$ and the $c_i$ are not identically
zero. Assume $m > m_1$. Then there is a polydisc
$\Delta_2 \subseteq \Delta_1$ about $0$ such that $c_1(z) \ne 0$ for
all $z \in \Delta_2^{(n-2)}$.
\end{lemma*}

\begin{proof}
Let $z$ be a fixed element of $\Delta_1^{(n-2)}$ and set
$g_z(z_{n-1}, z_n) = h(z, z_{n-1}, z_n)$. Regard $g_z$ as a polynomial
in $z_n$ over the field $\C\{z_{n-1}\}$ by replacing each coefficient
of $g_z$ by its power series expansion about $0$. Let $\zeta$ be a
primitive $m$-th root of unity and for $1 \le i \le m$, let
$\theta_i(z)$ denote the element
\[
c_1(z) \zeta^{(i-1) m_1} z_{n-1}^{m_1/m}
+ c_2(z) \zeta^{(i-1) m_2} z_{n-1}^{m_2/m} + \cdots
\]
of $\C\{z_{n-1}\}$. As the power series in $u$ obtained by expanding
\[
g_z(u^m, c_1(z) u^{m_1} + c_2(z) u^{m_2} + \cdots)
\]
is identically zero, we have, setting
$u = \zeta^{(i-1)} z_{n-1}^{1/m}$, that
$g_z(z_{n-1}, \theta_i(z)) = 0$. Thus $\theta_i(z)$ is a Puiseux root
of $g_z(z_{n-1}, z_n)$. But the $\theta_i(z)$, $1 \le i \le m$, are
distinct. (For suppose $1 \le i < j \le m$, and
$\zeta^{(i-1) m_k} = \zeta^{(j-1) m_k}$, $k = 1, 2, \ldots$. Let
$|u_1| < \rho_{n-1}$, $u_1 \ne 0$. Then the $m$ distinct roots of
$h(z, u_1^m, z_n)$ are the numbers
\[
\phi(z, u_1),\ \phi(z, \zeta u_1),\ \ldots,\ \phi(z, \zeta^{m-1} u_1).
\]
However,
\[
\begin{aligned}
\phi(z, \zeta^{i-1} u_1) &= c_1(z) \zeta^{(i-1) m_1} u_1^{m_1}
+ c_2(z) \zeta^{(i-1) m_2} u_1^{m_2} + \cdots \\
&= c_1(z) \zeta^{(j-1) m_1} u_1^{m_1}
+ c_2(z) \zeta^{(j-1) m_2} u_1^{m_2} + \cdots \\
&= \phi(z, \zeta^{j-1} u_1),
\end{aligned}
\]
a contradiction.) Hence the $\theta_i(z)$, $1 \le i \le m$,
constitute the complete set of Puiseux roots of
$g_z(z_{n-1}, z_n)$. For $1 \le i < j \le m$, let
$\eta_{i, j}(z) = \theta_i(z) - \theta_j(z)$. Then
\[
\eta_{i, j}(z) = d_1^{(i, j)}(z) z_{n-1}^{m_1/m}
+ d_2^{(i, j)}(z) z_{n-1}^{m_2/m} + \cdots,
\]
where
\[
d_k^{(i, j)}(z) = c_k(z) \zeta^{(i-1) m_k} - c_k(z) \zeta^{(j-1) m_k}
= c_k(z) \left( \zeta^{(i-1) m_k} - \zeta^{(j-1) m_k} \right).
\]
Let $1 \le i < j \le m$, and let $k_{i, j}$ be the least positive
integer $k$ such that $d_k^{(i, j)}(0) \ne 0$ (the number
$k_{i, j}$ certainly exists as $\theta_i(0) \ne \theta_j(0)$, and
hence $\eta_{i, j}(0) \ne 0$). Then
$O(\eta_{i, j}(0)) = m_{k_{i, j}} / m$. By continuity of the
functions $d_k^{(i, j)}$, there exists a polydisc
$\Delta_2 \subseteq \Delta_1$ about $0$ such that for every pair
$(i, j)$ with $1 \le i < j \le m$,
\[
d_{k_{i, j}}^{(i, j)}(z) \ne 0
\]
for every $z \in \Delta_2^{(n-2)}$. Thus, for all $(i, j)$,
\[
O(\eta_{i, j}(z)) \le m_{k_{i, j}} / m,
\]
for every $z$ in $\Delta_2^{(n-2)}$. We shall show that in fact we
have equality here.

For fixed $z$ in $\Delta_1^{(n-2)}$, let
$D_z(z_{n-1}) = F(z, z_{n-1})$ ($F$ the discriminant of $h$). Then as
$F(z, z_{n-1}) = z_{n-1}^r N(z, z_{n-1})$ by hypothesis,
\begin{equation}
\ord_0 D_z = r. \tag{4.2.3}
\end{equation}
On the other hand, $D_z(z_{n-1})$ can be regarded as the discriminant
of the polynomial $g_z(z_{n-1}, z_n) \in \C\{z_{n-1}\}[z_n]$ and
hence as an element of $\C\{z_{n-1}\}$. As
$\theta_1(z), \ldots, \theta_m(z)$ are the roots of
$g_z(z_{n-1}, z_n)$ in $\C\{z_{n-1}\}$, we have
\[
D_z = \prod_{1 \le i < j \le m} (\theta_i(z) - \theta_j(z))^2
= \prod_{1 \le i < j \le m} \eta_{i, j}(z)^2.
\]
Hence, by (4.2.3), we have
\begin{equation}
r = 2 \sum_{1 \le i < j \le m} O(\eta_{i, j}(z)). \tag{4.2.4}
\end{equation}
As (4.2.4) holds with $z = 0$, and as
$O(\eta_{i, j}(z)) \le O(\eta_{i, j}(0))$ for all $(i, j)$, it
follows that
\begin{equation}
O(\eta_{i, j}(z)) = O(\eta_{i, j}(0)), \tag{4.2.5}
\end{equation}
for all $(i, j)$.

Now there exists some $z'$ in $\Delta_2^{(n-2)}$ such that
$c_1(z') \ne 0$, as $c_1$ is not identically zero. Note that
$\zeta^{m_1} \ne 1$, as $0 < m_1 < m$ and $\zeta$ is a primitive
$m$-th root of $1$. Thus $\zeta^0 - \zeta^{m_1} \ne 0$ and hence by
(4.2.5), $O(\eta_{1, 2}(0)) = m_1/m$. Therefore, by (4.2.5),
$O(\eta_{1, 2}(z)) = m_1/m$, for all $z$ in $\Delta_2^{(n-2)}$: that
is, $d_1^{(1, 2)}(z) \ne 0$, for all $z \in \Delta_2^{(n-2)}$. As
\[
d_1^{(1, 2)}(z) = c_1(z) (\zeta^0 - \zeta^{m_1}),
\]
it follows that $c_1(z) \ne 0$ for all $z \in \Delta_2^{(n-2)}$.
This completes the proof of the lemma.
\end{proof}

\begin{lemma*}[4.2.8]
Let $h(z, z_n)$ be an irreducible Weierstrass polynomial of degree
$m \ge 1$ in the polydisc $\Delta_1$ about $0$ (in complex
$(n-1)$-space), and assume that $h$ satisfies the hypotheses of
Theorem~4.2.3. Let
\[
\begin{cases}
z_{n-1} = u^m \\
z_n = \phi(z, u)
\end{cases}
\]
be a parametrization for $h$, and let $m_1$ be the order of $\phi$
at $0$ in $u$. Then there is a polydisc $\Delta_2 \subseteq \Delta_1$
about $0$ such that
\[
\ord_{(z, 0, 0)} h = \min(m, m_1)
\]
for all $z \in \Delta_2^{(n-2)}$.
\end{lemma*}

\begin{proof}
Suppose first of all that $\phi \equiv 0$. Then $m = 1$ and
$h(z, z_n) = z_n$. Thus $\ord_{(z, 0, 0)} h = 1$ for all
$z \in \Delta_1^{(n-2)}$. Suppose on the other hand that $\phi$ is
not identically zero. Using the notation of Lemma~4.2.6 we may
assume, by refining $\Delta_1'$ suitably, that the power series
expansion for $\phi$ about $0$ is absolutely convergent in
$\Delta_1'$. Let this power series be arranged as follows:
\[
c_1(z) u^{m_1} + c_2(z) u^{m_2} + \cdots
\]
where $0 < m_1 < m_2 < \cdots$ and the $c_i$ are not identically
zero. Let $\C\{z\}$ be the field of all fractional power series in
$z_1, \ldots, z_{n-2}$. That is, inductively, $\C\{z\}$ consists of
all elements of the form
\[
a_h(z_1, \ldots, z_{n-3}) z_{n-2}^{h/m}
+ a_{h+1}(z_1, \ldots, z_{n-3}) z_{n-2}^{(h+1)/m} + \cdots
\]
where $m \ge 1$, $h$ is an integer (possibly negative), and
$a_k(z_1, \ldots, z_{n-3}) \in \C\{z_1, \ldots, z_{n-3}\}$. We can
show by induction on $n-2$, using Puiseux' theorem, that $\C\{z\}$
is algebraically closed. Let $K = \C\{z\}$, and regard
$h(z, z_{n-1}, z_n)$ as a polynomial in $z_n$ over the field
$K\{z_{n-1}\}$ by replacing each coefficient of $h$ by its power
series expansion about the origin (suitably arranged). Let $\zeta$
be a primitive $m$-th root of unity and for $1 \le i \le m$, let
$\theta_i$ be the element
\[
c_1(z) \zeta^{(i-1) m_1} z_{n-1}^{m_1/m}
+ c_2(z) \zeta^{(i-1) m_2} z_{n-1}^{m_2/m} + \cdots
\]
of $K\{z_{n-1}\}$. Now the power series in $z$ and $u$ obtained by
expanding
\[
h(z, u^m, c_1(z) u^{m_1} + c_2(z) u^{m_2} + \cdots)
\]
is identically zero. Hence, setting
$u = \zeta^{(i-1)} z_{n-1}^{1/m}$, we obtain
$h(z, z_{n-1}, \theta_i) = 0$. Thus $\theta_i$ is a Puiseux root of
$h(z, z_{n-1}, z_n)$. But the $\theta_i$ are distinct (cf.\ proof of
Lemma~4.2.7) and hence constitute the complete set of Puiseux roots
of $h(z, z_{n-1}, z_n)$. Thus
\[
h(z, z_{n-1}, z_n) = \prod_{i=1}^{m} (z_n - \theta_i)
\]
and so, where
\[
h(z, z_{n-1}, z_n) = z_n^m + a_1(z, z_{n-1}) z_n^{m-1}
+ \cdots + a_m(z, z_{n-1}),
\]
we have
\[
\begin{aligned}
a_1(z, z_{n-1}) &= -(\theta_1 + \cdots + \theta_m), \\
a_2(z, z_{n-1}) &= + \sum_{1 \le i < j \le m} \theta_i \theta_j, \\
&\vdots \\
a_m(z, z_{n-1}) &= (-1)^m \theta_1 \cdots \theta_m.
\end{aligned}
\]

We shall consider separately the two cases $m \le m_1$ and $m > m_1$.

\medskip
\noindent \textbf{Case I: $m \le m_1$.}
\medskip

We have $O(\theta_i) = m_1 / m \ge 1$, for $1 \le i \le m$. Hence
$O(a_j) \ge j$, for $1 \le j \le m$. This implies that, for
$1 \le j \le m$, the order (say $t_j$) in $z_{n-1}$ of the
holomorphic functions $a_j$ at the origin is at least $j$. Let
$\Delta_2 \subseteq \Delta_1$ be a polydisc about $0$ in which the
power series expansion about $0$ of each $a_j$ is absolutely
convergent. Then it follows that, for each point $z'$ in
$\Delta_2^{(n-2)}$, the order in $z_{n-1}$ of $a_j$ at $(z', 0)$ is
at least $j$. (This is clearly true if $a_j \equiv 0$, so assume
that $a_j$ is not identically zero, and let the power series
expansion for $a_j$ about $(0, 0)$ be arranged thus:
\[
a_{j, t_j}(z) z_{n-1}^{t_j} + a_{j, t_j + 1}(z) z_{n-1}^{t_j + 1} + \cdots
\]
where the $a_{j, k}$ are not identically zero. Then, where
$a_{j, k}'(z - z')$ denotes the power series expansion for $a_{j, k}$
about $z'$, we have that
\[
a_{j, t_j}'(z - z') z_{n-1}^{t_j}
+ a_{j, t_j + 1}'(z - z') z_{n-1}^{t_j + 1} + \cdots
\]
is the power series expansion for $a_j$ about $(z', 0)$. Thus the
order in $z_{n-1}$ of $a_j$ is $t_j$ ($\ge j$) at $(z', 0)$.)
Therefore, for each point $z'$ of $\Delta_2^{(n-2)}$, we have
$\ord_{(z', 0)} a_j \ge j$, hence
\[
\ord_{(z', 0, 0)} (a_j(z, z_{n-1}) z_n^{m-j}) \ge j + m - j = m.
\]
It follows that
\[
\ord_{(z', 0, 0)} h = m = \min(m, m_1)
\]
for every $z' \in \Delta_2^{(n-2)}$.

\medskip
\noindent \textbf{Case II: $m > m_1$.}
\medskip

By Lemma~4.2.7, there is a polydisc $\Delta_2 \subseteq \Delta_1$
about $0$ such that $c_1(z) \ne 0$ for all
$z \in \Delta_2^{(n-2)}$. As in Case~I, $O(\theta_i) = m_1 / m$.
Thus $O(a_j) \ge j m_1 / m$, for $1 \le j \le m$. Therefore, for
$1 \le j \le m$, the order in $z_{n-1}$ of the holomorphic function
$a_j$ at the origin is greater than or equal to $j m_1 / m$. Assume,
by refining $\Delta_2$ if necessary, that the power series expansion
about $0$ of each $a_j$ is absolutely convergent. Then it follows
that, for $1 \le j \le m$, and for each point $z' \in
\Delta_2^{(n-2)}$, the order in $z_{n-1}$ of $a_j$ at $(z', 0)$ is
greater than or equal to $j m_1 / m$. Therefore, for $1 \le j < m$
and for $z' \in \Delta_2^{(n-2)}$, we have
$\ord_{(z', 0)} a_j \ge j m_1 / m$, hence
\[
\begin{aligned}
\ord_{(z', 0, 0)} (a_j(z, z_{n-1}) z_n^{m-j})
&\ge j m_1 / m + m - j \\
&= m - j (1 - m_1 / m) \\
&> m - m (1 - m_1 / m) \\
&= m_1.
\end{aligned}
\]
Now $a_m(z, z_{n-1}) = (-1)^m \theta_1 \cdots \theta_m$ and
\[
\theta_1 \cdots \theta_m = d_{m_1}(z) z_{n-1}^{m_1}
+ d_{m_1 + 1}(z) z_{n-1}^{m_1 + 1} + \cdots,
\]
where $d_{m_1}(z) = (-1)^{m-1} c_1(z)^m$. Thus, as
$c_1(z') \ne 0$ for each $z' \in \Delta_2^{(n-2)}$, we have that
$\ord_{(z', 0)} a_m = m_1$ for each $z' \in \Delta_2^{(n-2)}$.
Hence
\[
\ord_{(z', 0, 0)} h = m_1 = \min(m, m_1)
\]
for every $z' \in \Delta_2^{(n-2)}$.
\end{proof}

We can give at last the

\begin{proof}[Proof of Theorem 4.2.3]
Factor $h$ into a product of irreducible Weierstrass polynomials
$h_1 \cdots h_k$. Let $1 \le i \le k$ and let $G(z)$ be the
discriminant of $h_i(z, z_n)$. Recall that the discriminant $F(z)$
of $h(z, z_n)$ has the form
\[
F(z) = z_{n-1}^r N(z),
\]
where $N(z)$ is holomorphic and non-vanishing in $\Delta_1$. Let
$\delta = (\delta_1, \ldots, \delta_{n-1})$ be the polyradius of
$\Delta_1$. Now $F(z) = G(z) H(z)$, for some holomorphic $H(z)$, by
Theorem~2.3.3. Hence, for every fixed $z \in \Delta_1^{(n-2)}$, we
must have
\[
G(z, z_{n-1}) = z_{n-1}^s \cdot U(z, z_{n-1})
\]
for some $s$, $0 \le s \le r$, and some $U(z, z_{n-1})$, holomorphic
and non-vanishing in the disc $|z_{n-1}| < \delta_{n-1}$. By zero
system continuity, $s$ does not depend on $z$. In fact,
$U(z, z_{n-1})$ is holomorphic in the variables
$z_1, \ldots, z_{n-2}$, as well as $z_{n-1}$, as the expression
\[
U(z, z_{n-1}) = \frac{1}{2 \pi i} \int_C \frac{U(z, t) \, dt}
{t - z_{n-1}}
\]
holds for $z \in \Delta_1^{(n-2)}$ and $z_{n-1}$ inside $C$, for any
circle $C$ about $0$ in the disc $|z_{n-1}| < \delta_{n-1}$. Thus
$h_i$ satisfies the hypotheses of Theorem~4.2.3. Let $p$ be the
degree of $h_i$. By Lemma~4.2.6 there exists a parametrization
\[
\begin{cases}
z_{n-1} = u^p \\
z_n = \phi_i(z, u)
\end{cases}
\]
for $h_i$ in $\Delta_1$. Let $p_1$ be the order of $\phi_i$ at $0$
in $u$. By Lemma~4.2.8 there exists a polydisc
$\Delta_2^{(i)} \subseteq \Delta_1$ about $0$ such that
\[
\ord_{(z', 0, 0)} h_i = \min(p, p_1)
\]
for all $z'$ in $\Delta_2^{(i)(n-2)}$. Thus $h_i$ is order-invariant
in $\Delta_2^{(i)(n-2)} \times \{ (0, 0) \}$. Let
$\Delta_2 = \bigcap_{i=1}^{k} \Delta_2^{(i)}$. Then
$h \left( = \prod_{i=1}^{k} h_i \right)$ is order-invariant in
$\Delta_2^{(n-2)} \times \{ (0, 0) \}$.
\end{proof}

\end{document}
